\subsection{Further investigation of the original frontier}
The previous release left $34$ exact-M3 rings unresolved. This continuation supplies $25$ further general complex exclusions, leaving $9$ questions. The earlier ledger and replay receipts are retained byte for byte; accepted changes are recorded in a separate overlay ledger.
\subsection{A simultaneous fixed-point module obstruction}
\label{subsec:simultaneous-fixed-point}

\begin{theorem}
The fusion rings $\sha{R8_3_1203}$, $\sha{R8_3_1350}$, $\sha{R8_3_2842}$, and
$\sha{R8_3_2885}$ have no categorification over $\mathbb C$.
\end{theorem}
\begin{proof}
We give a general reduction and then its exact finite verification.
Let $g$ be an invertible object of order two fixing a simple object in a
categorification of a commutative fusion ring.  A fixed simple gives a
rank-one module category over $\operatorname{Vec}(C_2,\omega)$, so its
module associativity trivializes $[\omega]$.  Thus $A=\mathbf1\oplus g$
is a connected separable algebra.  In the category $\mathcal C_A$ of
right modules, a two-element $g$-orbit induces a single simple module,
whereas a fixed simple $V$ induces two simples $V_+,V_-$, each with
underlying object $V$.  These statements follow from induction
adjunction and semisimplicity.

The span of induced objects is invariant under every action matrix
$M_X$, and under its transpose $M_{X^*}$ by rigidity.  Its orthogonal
complement, with basis $V_+-V_-$ indexed by the fixed simples, therefore
carries matrices $S_X$ satisfying
\[
 S_{\mathbf1}=I,\qquad S_{X^*}=S_X^t,\qquad
 S_XS_Y=\sum_Z N_{XY}^Z S_Z.
\]
The corresponding two-by-two module block has both row and column sums
$a=N_{XV}^W$; hence it is
$\left(\begin{smallmatrix}u&a-u\\a-u&u\end{smallmatrix}\right)$.
Consequently $(S_X)_{VW}\in\{-a,-a+2,\ldots,a\}$.
The matrix $S_g$ is diagonal with entries $\pm1$.
If $gX=X$ and $N_{XV}^V$ is odd, then
$S_gS_X=S_X$ forces $(S_g)_{VV}=1$.  For the four rings in the theorem,
every fixed vertex admits such an $X$, and therefore $S_g=I$.

The fixed indices are $4,5,6,7$, with $g$ at index $3$ in the first two
rings and index $1$ in the last two.  The representation $S$ factors
through $R_{\mathbb Q}/(g-1)$, a six-dimensional rational algebra
generated by the self-dual basis object of index $4$.  The annihilating
polynomial for its image is
\[
\begin{cases}
(t^2-14t+4)(t^4+2t^3-6t^2-12t-4),&1203,2842,\\
(t^2-12t-14)(t^4-10t^2+12t-2),&1350,2885.
\end{cases}
\]
Every $S_X$ is a specified rational polynomial in $S_4$.  The exact
polynomial identities are checked independently in the original
rank-eight regular representation by multiplication by $I+N_g$.

We enumerate all symmetric integral $4\times4$ matrices $S_4$ having
the prescribed bounds and parities, test the displayed polynomial,
and reconstruct the other $S_X$.  A simultaneous bounded integral
representation never occurs:
\[
\begin{array}{c|r|r|r}
\text{ring}&\text{all matrices}&\text{polynomial survivors}&
\text{simultaneous survivors}\\\hline
1203&110592&32&0\\
1350&62208&16&0\\
2842&110592&32&0\\
2885&62208&16&0
\end{array}
\]
All computations use exact rational or unbounded integer arithmetic;
the finite bounds come from the module-block argument.  The discovery
enumerator uses a simultaneous diagonal-sign gauge and is replayed by
an independent enumerator without this reduction.  Both retain their
checks under optimized Python.  As a positive control, the known
model $R8\_3\_342=\operatorname{Rep}(C_{13}\rtimes C_6)$ passes every
simultaneous multiplication, duality, bound and parity check on its
two-dimensional difference representation.  The evidence and two
read-only verifiers are in
\path{evidence/rank8_frontier_20261021/quadratic/}.
\end{proof}

The argument uses only semisimple module-category theory (see Ostrik,
\emph{Module categories, weak Hopf algebras and modular invariants},
Transformation Groups \textbf{8} (2003), 177--206, and
Etingof--Gelaki--Nikshych--Ostrik, \emph{Tensor Categories}, Chapter~7).
In particular, neither pivotality nor unitarity is assumed.  The
conclusions for $1350$ and $2885$ strengthen the previous unitary-only
obstructions.

\begin{theorem}
The fusion ring $\sha{R8_3_921}$ has no categorification over $\mathbb C$.
\end{theorem}
\begin{proof}
Use the same split module, now with $g$ at index $5$ and fixed indices
$3,4,6,7$.  Put $B=S_3$.  Since index $4$ is dual to $3$ and the ring
is commutative, $BB^t=B^tB$.  Thus $B$ is a real normal matrix.
All sixteen fixed-block coefficients of index $3$ are odd, and its four
diagonal coefficients are $1$.  Hence
$|\operatorname{tr}B|\leq4$ and $\|B\|_F^2\geq16$.
Its squarefree source annihilator is
\[
 p(t)=t(t^2-8t+6)(t^2+2t+4).
\]
The characteristic polynomial of this rational four-dimensional matrix
is $t^a(t^2-8t+6)^b(t^2+2t+4)^c$, where
$a+2b+2c=4$.  Its trace is $8b-2c$, and normality gives squared
Frobenius norm $52b+8c$.  The trace bound rules out $b>0$; the norm
bound then forces $c=2$.  Thus $q(B)=0$, where $q(t)=t^2+2t+4$.
The exact identity in the regular fusion representation
\[
 q(N_3)(-7N_1-7N_2-3N_3+30N_6-7N_7)=148N_6
\]
therefore forces $S_6=0$.  This contradicts the odd fixed-block entry
$N_{6,3}^{3}=1$.  The argument needs no condition $S_g=I$.
The verifier \path{normal_split_spectral.py} checks all source
identities and all six characteristic-factor possibilities.  Its
independent full-box search examines all $131072$ bounded matrices:
$2096$ are normal and $16$ satisfy $p(B)=0$, and every one of these
satisfies $q(B)=0$.  The matrix norm here is an ordinary exact
linear-algebra invariant; it introduces no categorical unitarity
assumption.
\end{proof}

\subsection{A general inversion obstruction from the $C_3+3$ TPE}
\begin{proposition}
Every complex near-group category of type $(C_3,3)$ is Galois conjugate to
a unitary category. Consequently $\sha{R8_3_436}$ has no complex categorification.
\end{proposition}
\begin{proof}
Thornton's sphericality and Galois-pseudo-unitarity theorem
\cite[Theorem~1.3]{ThorntonNear} allows us to take the noninvertible
dimension $d=(3+\sqrt{21})/2$. Its second indicator is $+1$ by
\cite[Theorem~3.9]{LPRV}, since its self-multiplicity is odd.
We may therefore use the complete spherical tetrahedral parametrization
of \cite[Theorem~2.3, Proposition~2.7, Theorem~4.12]{NearGroupTPE2026}.

Write $w^2+w+1=0$, $q=1+w$, and $t=-1/d=(3-d)/3$.
Both possible nondegenerate bicharacters are $\chi(g,h)=w^{gh}$;
the even refinement is forced to be $a(0)=1$, $a(1)=a(2)=w$.
Put $x=\beta(1)$ and $y=\beta(2)$.
The product identity and two cubic equations give
\[
xy=w/3,\qquad y^2+qtx=tw,\qquad x^2+qty=tw.
\]
Thus $(y-x)(x+y-qt)=0$. The case $x=y$ would imply
$6t^2-6t+1=0$, contrary to $t^2-t=1/3$.
Hence $x+y=qt$, $xy=w/3$, and the rotation scalar is uniquely
$c=1/(2+w)$. The exact verifier checks that these relations are equivalent
to all nine cubic equations, including both dimension roots and both
bicharacters.

The numbers $u=x/q$ and $v=y/q$ are roots of the real polynomial
$z^2-tz+1/3$, whose discriminant is $-d/3<0$.
Therefore $v=\overline u$, so
$\overline x=y/w$, $\overline y=x/w$.
These are the Hermitian scalar conditions, and
\cite[Corollary~5.5]{NearGroupTPE2026} unitarizes the complete table of
the original category. The completeness of that table is essential;
an existence-only reconstruction would not justify this conclusion.

For the source tensor, the indices $[0,1,2,5]$ give this near-group ring,
and the invertible at index $3$ conjugates indices $1$ and $2$.
Galois conjugation transports this ordinary tensor autoequivalence along
with the category, preserving its label permutation. The unitary
classification \cite[Theorem~9.1, Example~9.4]{IzumiNear} identifies
the resulting neutral category with an Izumi--Xu category.
Both have trivial outer group of ordinary complex tensor
autoequivalences \cite[Section~4.2, p.~10]{GJSCyclic}; inner
autoequivalences fix labels since the ring is commutative.
The required inversion is impossible.
\end{proof}
This new argument uses the full spherical parametrization in the October
2026 preprint. The exact checker verifies all source coefficients and the
polynomial reduction; the cited categorical inputs and the new argument
remain subject to mathematical review.

\subsection{A uniform $\operatorname{Rep}(A_5)$ complement obstruction}
\label{sec:rank8-a5-complement}

The integral ring $\mathtt{R8\_3\_3344}$ is excluded over $\mathbb C$.
The argument separates a structural reduction from a complete relative
cocycle calculation; neither step assumes unitarity.

\paragraph{Structural reduction.}
Suppose a fusion category $\mathcal C$ of dimension $360$ contains
$\mathcal D\simeq\operatorname{Rep}(A_5)$, with three complementary
simples of dimension $10$, on which the left and right action matrices
of the simples of dimensions $1,3,3,4,5$ are
\[
 I,\quad J,\quad J,\quad I+J,\quad 2J-I,
\]
where $J$ is the $3\times3$ all-ones matrix. Then $\mathcal C$ is
group-theoretical.
Indeed, the complement $\mathcal M$ is an indecomposable
$\operatorname{Rep}(A_5\times A_5)$-module category. By Ostrik's
classification it is $\operatorname{Rep}_{\alpha}(K)$ for a subgroup
$K\leq A_5\times A_5$. The displayed action forces all three projective
irreducible degrees to equal an integer $q$, hence $|K|=3q^2$.
Neither projection of $K$ can be $A_5$: Goursat's lemma reduces this
case to a graph or a direct product, and every projective sector of
$A_5$ has rank at least four. The spectrum $\{3,0,0\}$ of the degree-three
action forces each projection to contain a $3$-cycle. Thus its proper
projections have orders $3,6,12$, leaving $|K|=3,12,48$.

Set $A=\operatorname{Fun}(A_5)$ and
$\mathcal E={}_A\mathcal C_A$. Its diagonal subcategory is
$\operatorname{Vec}_{A_5}$. The complementary $A$-bimodules are
projectively $K$-equivariant vector bundles on the free $K$-set
$A_5\times A_5$. They number $3600/|K|$ and form a transitive
$A_5\times A_5$-set. Their common Frobenius--Perron dimension is therefore
$q/2$, since their squared dimensions sum to $300$.
The possibilities are
\[
 (|K|,q,\operatorname{rank},\operatorname{FPdim})
 =(3,1,1200,\tfrac12),\quad(12,2,300,1),\quad(48,4,75,2).
\]
The first is impossible and the second makes $\mathcal E$ pointed.
In the third, Goursat's lemma gives
$K=A_4\times_{C_3}A_4$. Each complementary simple has left stabilizer
$V_4$, so $XX^*$ is the sum of its four stabilizing invertibles.
Consequently $\mathcal E_{\rm ad}=\operatorname{Vec}_{A_5}$, since the
conjugates of $V_4$ generate $A_5$.
Each universal-grading component is an invertible
$\operatorname{Vec}_{A_5}$-bimodule. But a component with left stabilizer
$V_4$ has dual $\mathcal C(A_5,1,V_4,\psi)$, which is not pointed:
distinct Sylow $2$-subgroups of $A_5$ intersect trivially, giving a
simple of dimension $[V_4:1]=4$ by the double-coset formula.
This contradicts invertibility and proves the reduction.

The unique proper subring of $\mathtt{R8\_3\_3344}$ has basis
$[0,1,2,6,7]$ and the $\operatorname{Rep}(A_5)$ tensor.
Every categorification of this simple dimension-$60$ subring is
$\operatorname{Rep}(A_5)$ by Etingof--Nikshych--Ostrik,
\emph{Weakly group-theoretical and solvable fusion categories},
Theorem~9.12. The full tensor verifies the hypotheses above.

\paragraph{All group-theoretical possibilities.}
The ring has no nontrivial invertible and trivial universal grading,
so its center has no nontrivial invertible. In a presentation
$\mathcal C(L,\omega,H,\psi)$, this forces $L$ to be perfect.
An exhaustive check of the $162$ groups of order $360$ leaves only
$L=A_6$. The identity double-coset sector $\operatorname{Rep}(H)$
forces $H$ perfect as well. The perfect-group catalogue on the divisors
of $360$, together with the index-three action excluding order $120$,
leaves $H=1,A_5,A_6$. The endpoints have the wrong ranks. The index-six
coset action identifies every $A_5$ embedding with the usual point
stabilizer up to an automorphism of $A_6$.

The restriction $H^2(A_6,\mathbb C^\times)\to
H^2(A_5,\mathbb C^\times)$ is surjective, and
\[
 \ker\bigl(H^3(A_6,\mathbb C^\times)\longrightarrow
 H^3(A_5,\mathbb C^\times)\bigr)\simeq C_2.
\]
Here is a reproducible calculation rather than an assumed cohomology
table. Normalized mod-$2$ bar cohomology on a Sylow $D_8$, with all
$360$ stable-element conjugation conditions, gives $2$ stable classes
in degree $2$ and $4$ in degree $3$. The restriction kernels on the
point-fixing $V_4$ have orders $1$ and $2$, respectively.
The degree-two result and the coefficient exact sequence prove the
surjectivity above and bound the relative $2$-part by $C_2$.
Transfer annihilates the relative kernel by the index $6$.
For its $3$-part, restriction to the Sylow $C_3^2$ embeds it in
$H^4(C_3^2,\mathbb Z)=\langle x^2,xy,y^2\rangle_{\mathbb F_3}$.
An explicit normalizer rotation leaves only the line
$\langle x^2+y^2\rangle$, which restricts injectively to the coordinate
$C_3$ in $A_5$. Thus the relative $3$-part vanishes.
The previously certified relative $\mu_2$ cocycle has a nontrivial
$A_4$ stabilizer multiplier and supplies the nonzero relative class.

There are consequently only two relative gauge classes of
$(\omega,\psi)$. The trivial one gives rank $9$.
The nontrivial one is the already verified model
$\mathtt{R8\_3\_3342}$. Its three dimension-ten simples satisfy
$N_{X_iX_i}^{X_i}=2$, whereas all three corresponding coefficients
of $\mathtt{R8\_3\_3344}$ equal $3$. Hence no group-theoretical model
exists for the latter ring. The structural reduction proves its
noncategorifiability over $\mathbb C$.

\paragraph{A general obstruction for the dimension-$660$ ring.}
The ring $\mathtt{R8\_3\_731}$ is also excluded over $\mathbb C$.
Its source dimensions are $1,5,5,10,10,11,12,12$ and its based ring is
simple. Any categorification is integral, hence pseudo-unitary and
spherical. It cannot be braided: simplicity and the prime-dimensional
simple of dimension $5$ exclude the nondegenerate case by
Etingof--Nikshych--Ostrik Corollary~7.2, while the symmetric case would
be a group representation category. The exhaustive ordinary-group
scan treats all $40$ groups of order $660$ and finds no matching tensor.
Liu--Palcoux--Ren, \emph{Interpolated family of non group-like simple
integral fusion rings of Lie type}, Theorem~5.1, therefore forces its
hypothetical Drinfeld center to be simple. Its dimensions are integral,
so Corollary~7.2 excludes every nontrivial central simple of prime-power
dimension. This strategy is already proposed in their Remark~5.15;
the new result is the complete finite induction and rotation certificate
below.

The formal codegrees are $660,4,5,5,11,11,12,12$, so the distinguished
central simple dimensions are $1,165,132,132,60,60,55,55$.
The exact Gram ellipsoid gives $38058$ vectors, whose seven required
unit-multiplicity dimension slices have sizes
$65,61,61,31,31,37,37$. Complete exact row-sum and positive-semidefinite
completion leaves precisely $450$ low induction frames. The required
duality invariance of the restriction-column multiset excludes $434$.
For each remaining frame, all admissible high restriction columns are
enumerated exactly, with central dimension dividing $660$ and not a
prime power. Exact Gram-completion Farkas certificates exclude another
$13$ frames.

The last three frames, with indices $293,301,303$, have respectively
$280,210,324$ admissible high columns. Their second indicators are all
forced, namely $(0,0,0,0,1,1,1)$ in nonunit source order.
Joint Gram, twist-moment, and tensor-rotation equations through order
six use $53$ complete normalized rational trace patterns. Divisors of
$19800$ together with the zero pattern exhaust these patterns; this is
a trace enumeration bound, not a claimed exponent bound. The three
integer systems have $182$ equations each and respectively
$14917,11207,17249$ variables. Each has a saved integer Farkas vector
$y$ satisfying $A^{\mathsf t}y\geq0$ and $b^{\mathsf t}y<0$.
The read-only verifier reconstructs all $450$ frames and all three
integer systems and checks every certificate exactly, without calling
an LP solver. Thus $434+13+3=450$ branches are excluded, proving general
noncategorifiability. The previous unitary-only obstruction is now
supplemented by this independent general obstruction. The printed
$F_{660}$ in the triangular-prism paper remains
$\mathtt{R8\_3\_722}$, a different tensor.

\paragraph{Evidence.}
The full proof, exact source hashes, all subgroup and cohomology
receipts, and normal/optimized Python replays are in
\path{evidence/rank8_frontier_20261021/integral/}.
The independent review recomputes the mod-$2$ calculation without
importing the discovery verifier. The categorical inputs are Ostrik's
module-category classification (\S3.4, Theorem~2),
Etingof--Nikshych--Ostrik's invertibility of graded components
(\emph{Fusion categories and homotopy theory}, Theorem~6.1),
Gelaki--Naidu's double-coset simple-object formula (Theorem~5.1),
and Cartan--Eilenberg stable elements and restriction/transfer.

\subsection{Exact structural and infinite-family construction filters}
The following results concern specified construction domains; they are not
general categorification exclusions for the rings that remain open.
The adjoint closure is the full ring for
$\sha{R8_3_65}$, $\sha{R8_3_439}$, $\sha{R8_3_567}$,
$\sha{R8_3_901}$ and $\sha{R8_3_1463}$, excluding a nontrivial faithful group
grading of these particular categories. The two latter rings have a unique
Fibonacci generator. In a putative Fibonacci Deligne factorization, the
complementary simple basis would be forced to be $[0,1,2,7]$ and
$[0,1,3,6]$, respectively; neither basis is closed under multiplication.
They also cannot be braided: the Fibonacci subcategory would be
nondegenerate, since a symmetric fusion category is integral, and would
split off by \cite[Corollary~2.46 and Theorem~3.13]{DGNO}.
We do not infer a coherent central lift from based-ring commutativity.
The generalized near-group $\sha{R8_3_65}$ has a free $C_4$ action on its four
noninvertibles and noninvertible FP dimension satisfying $d^2=1+11d$.
The classification \cite[Theorem~1]{SchopierayBraidedGNG} therefore excludes
a braided categorification as well.

There is a stronger unbraided construction restriction for $\sha{R8_3_65}$.
Any canonical $\operatorname{Rep}(G)$ of an equivariantization must lie in
its pointed $C_4$ subcategory, since its other simple dimensions are
irrational. Hence every nontrivial possibility contains a central
$\operatorname{Rep}(C_2)$. Its sign object acts freely on all eight source
simples. The orbit description of equivariantization forces a rank-four
de-equivariantization of global dimension $125+55\sqrt5>24+12\sqrt3$,
contrary to the unrestricted rank-four bound proved in the rank-four
companion, Corollary~\path{finite:cor:sharp-bounds}. Thus every
nontrivial finite-group equivariantization is excluded for this ring.
In contrast, the complete rank-two based-module equations have exactly
two solutions up to exchanging vertices:
\[
N_\rho=\begin{pmatrix}0&1\\1&11\end{pmatrix},\qquad
N_\rho=\begin{pmatrix}3&5\\5&8\end{pmatrix}.
\]
An actual module category with either action would trivialize the pointed
$H^3(C_4,\mathbb C^\times)$ class. Neither that vanishing nor an
associator for the second action is inferred from these matrices.

For $\sha{R8_3_439}$ and $\sha{R8_3_567}$, the complete $C_3$ split-module
calculation permits exactly two conjugate nontrivial Fourier systems.
The subgroup algebra $A=\bigoplus_{g\in C_3}g$ exists because a fixed
noninvertible trivializes its pointed third-cohomology class. Its Morita
dual has six invertibles and two noninvertibles. The right deck actions
and the full left fusion action force the dual noninvertible matrix via
an integer linear system of rank $64$ in $64$ unknowns. The eight dual
actions are independent, and all $512$ resulting coefficients recover
the original respective source ring. This closes this particular Morita
route at the level of based rings; it asserts neither equivalence of
associators nor existence of a categorification.

We also enumerate every possible underlying object of a connected
separable algebra in seventeen source rings, supplementing the six
quadratic cases treated separately.
The coefficient bounds $a_0=1$, $a_i=a_{i^*}$ and
$0\leq a_i\leq\lfloor\operatorname{FPdim}(X_i)\rfloor$, invertible
stabilization, and the entrywise nonnegative positive-semidefinite
induction residual Gram condition leave the following numbers of objects:
\[
\begin{array}{c|rrrrrrrrr}
\text{source index}&65&439&455&567&721&901&903&928&941\\
\text{survivors}&30&28&21&22&45&18&27&24&38\\[2pt]
\text{source index}&982&985&1154&1463&1615&1617&1660&1661&\\
\text{survivors}&31&31&81&26&49&39&45&49&
\end{array}
\]
Every retained object is independently checked with rational arithmetic,
and ten complete finite domains are independently re-enumerated in
Python; the remaining domains use the same exactly bounded native
enumerator. For the noncommutative tensors
the induction matrix uses right multiplication by the candidate algebra.
No surviving object is asserted to carry an associative multiplication.
For $\sha{R8_3_65}$, every object $1+X$ with $X$ noninvertible fails this
necessary filter, leaving only the order-two pointed choice among
objects of the form $1+X$.

For the new cyclic family $\mathrm{NG}(C_n,n)$ of
\cite[Theorem~2.10 and Remark~2.11]{NearGroupTPE2026}, and indeed for every
categorification of those rings, the FP global dimension $D_n$ has
\[
\operatorname{Tr}(D_n)=n(n+4),\qquad
\operatorname{Nm}(D_n)=n^2(n+4).
\]
If a target quadratic FP dimension $D$ equals $sD_n$ with $s\in\mathbb Q_{>0}$,
then necessarily
\[
n=\frac{\operatorname{Tr}(D)^2}{\operatorname{Nm}(D)}-4,
\qquad
s=\frac{\operatorname{Tr}(D)}{n(n+4)}.
\]
Thus an exact calculation replaces a bounded search over $n$. Direct and
Morita equivalence require $s=1$; a group extension or equivariantization
requires integral $s$. If the target is a de-equivariantization by a group
of order $k=1/s$, then $k\mid n$, its rank is $n/k+k$, and its invertible
group is cyclic of order $n/k$. To see this, the integral subcategory
$\operatorname{Rep}(G)$ in the near-group category must be pointed, so
$G$ is cyclic. All its stabilizers have trivial second cohomology.
The equivariantization orbit description
\cite[Corollary~2.13 and Remark~3.1]{BurciuNataleEquivariantization}
then forces every invertible orbit to be fixed and the unique
noninvertible orbit to be free. Finally, $\mathrm{NG}(C_n,n)$ has trivial
universal grading, excluding proper neutral-component extraction.

These exact dimension, rank and invertible-group tests, supplemented by
the independent general exclusion of $\sha{R8_3_436}$, exclude all five
single-operation routes for the original $34$ frontier tensors, for every
positive $n$. For example, $\sha{R8_3_65}$ forces $(n,s)=(121,4/121)$;
$\sha{R8_3_439}$ and $\sha{R8_3_567}$ force $(n,k)=(12,2)$, whose
de-equivariantization would have cyclic, rather than $S_3$, invertibles.
Arbitrary chains of operations, Morita changes inserted into those chains,
and general noncentral algebra condensations are outside this claim.

For $\sha{R8_3_455}$, a separate exact Fourier calculation excludes the
self-Morita-duality required of the neutral category within the strict
Evans--Gannon $\mathrm{HI}(C_3)$ matrix presentation, including its known
Galois companions. We retain the gap between an arbitrary abstract complex
neutral category with vanishing pointed cocycle obstructions and that
strict presentation. This conditional calculation does not change the
general status of $\sha{R8_3_455}$.

\subsection{Closing the remaining near-group and quantum-group model branches}
\label{subsec:frontier-known-family-closure}
There is a useful Morita-invariant construction obstruction.
If a category has no nontrivial invertible basis object commuting with
all simples and has full adjoint subcategory, its center has no
nontrivial invertibles. Indeed a central invertible must forget to the
unit; its half-braiding is a tensor automorphism of the identity and is
trivial on every constituent of $X\otimes X^*$, hence on the full
adjoint closure. Both $\sha{R8_3_439}$ and $\sha{R8_3_567}$ satisfy these
hypotheses exactly. Morita equivalence preserves the braided center
\cite[Introduction]{ENOExt}.

The complete $SU(2)$ product-dimension sieve retains just the branch
$2A_{12}$ for these two rings: one adjoint $SU(2)_{10}$ factor and total
integer multiplier two. Full $SU(2)_{10}$ has a faithful $C_2$ grading;
each other allowed realization of this dimension branch uses precisely
one dimension-two extension, equivariantization or Deligne factor.
Each supplies a nontrivial central invertible, preserved under any
Morita changes. Thus the whole retained branch is excluded, including
all ADE module-category Morita duals of $SU(2)_{10}$. Together with the
thirty dimension obstructions this gives thirty-two $SU(2)$ family
exclusions among the original thirty-four rings. The remaining two
dimension matches are integral rings already excluded independently.
No assertion about arbitrary algebra condensation or other Lie types
is made.

The actual inversion action of $C_2$ on the known
$\mathrm{NG}(C_3,6)$ also deserves an explicit check. At the normalized
solution, \cite[Lemma~13.1 and Example~13.10]{IzumiNear} gives the
involution $(-\mathrm{id},R_{\pi/4}L)$, with $L=\operatorname{diag}(1,-1)$
and $(R_{\pi/4}L)^2=I$. Its crossed product is coherent, but its tensor
has multiplicity six and a nontrivial $C_2$ grading, so it cannot be
either source tensor. Under the additional hypothesis that the full
pointed $S_3$ associator class is trivial, condensation of
$\bigoplus_{g\in S_3}g$ instead forces a rank-five dual with dimensions
$1,1,2,3+2\sqrt3,3+2\sqrt3$. Complete associativity and rigidity leave
five based tensors. Full coefficient matching and fresh replay of four
frozen rank-five exclusion certificates leave only $R5\_3\_40$.
The known equivariantization realizing that based ring has a central
$\operatorname{Rep}(C_2)$ and so cannot be this dual. This structural filter alone does not decide other associators; the centerless Morita obstruction in Section~\ref{subsec:r439-centerless-morita} excludes the required dual.

The new indicator calculation now decides which pointed-cocycle branch
is necessary. After the proved Galois-pseudo-unitary reduction, the
complete joint-rotation test on the unique full-character-$S$ low frame
leaves only
\[
 (\nu_2(X_1),\ldots,\nu_2(X_7))=
 \begin{cases}
 (0,0,1,1,1,1,1),&R8\_3\_439,\\
 (0,0,0,0,-1,-1,-1),&R8\_3\_567.
 \end{cases}
\]
All fourteen other indicator branches have exact separating witnesses.
For an order-two invertible in the canonical positive spherical
structure, its indicator is the sign of its pointed $H^3(C_2,\mathbb
C^\times)$ class; see \cite[Definition~2.2 and Proposition~4.1]{SchauenburgPointedFrontier}
with trivial subgroup. The $C_3$ restriction is already trivialized by
a fixed noninvertible. Restriction followed by corestriction is
multiplication by the subgroup index for arbitrary subgroups
\cite[Section~3.3, equation~(3.12)]{BrownTransferFrontier}; hence the
two trivial Sylow restrictions imply $2\omega=3\omega=0$ and
$\omega=0$. Galois transport preserves vanishing and transports the
subgroup algebra. Thus every complex categorification of $\sha{R8_3_439}$
has the rank-five Morita dual just described, necessarily with a tensor
category structure inequivalent to the known equivariantization.
For $\sha{R8_3_567}$, the $C_2$ restriction is nontrivial and the full
$S_3$ subgroup algebra is impossible. Its $C_3$ condensation remains
available. Neither conclusion proves existence or nonexistence.

Finally let $D_0$ be the irrational global FP dimension of an arbitrary
near-group category with group order $n$ and multiplicity $m$. Then
$\operatorname{Tr}(D_0)=4n+m^2$ and
$\operatorname{Nm}(D_0)=n(4n+m^2)$. For $\sha{R8_3_721}$,
$D=136+32\sqrt{17}$ has trace $272$ and norm $1088$.
An equality $D=tD_0$ with positive integer $t$ therefore forces
$n=4/t$ and $m^2=256/t$. The complete possibilities are
$(n,m,t)=(4,16,1)$ and $(1,8,4)$; $t=2$ gives a nonsquare.
Both groups of order four are excluded by full tensor matches to
$R5\_16\_1909$ and $R5\_16\_1653$ and fresh frozen
certificate replay. The rank-two case is excluded by
\cite[Main Theorem]{Ostrik2}.
This exhausts all near-group orders and multiplicities with integer
dimension multipliers and Morita changes for this target, but not
dimension-decreasing operations or arbitrary nonintegral products.

\subsection{Complete quadratic module and construction domains}

For the six residual indices $1055,1768,1795,3258,3413,3472$, connected
rank-two based modules admit a complete finite classification. The
positive Frobenius--Perron character and its quadratic conjugate exhaust
their two-dimensional rational representation. If a self-dual generator
has dimension polynomial $t^2-Tt+U$, its action is
\[
 \begin{pmatrix}a&b\\b&T-a\end{pmatrix},\qquad
 0\leq a\leq T,\qquad b^2=a(T-a)-U>0.
\]
Vertex interchange allows $a\geq T-a$. All other action matrices follow
from rational affine expressions in the generator dimension. Exact
nonnegative integrality, duality, and all fusion products leave respectively
$1,1,2,1,1,2$ based-module classes. Their full actions and candidate
internal-end objects are recorded. These data do not solve module
associativity and are not constructions of categories.

For index 1055 the forced $1+g$-module difference representation has rank
two. Normality and the source annihilator reduce the first generator to
zero; all 1296 remaining bounded simultaneous possibilities give exactly
four representations, each with $S_g=-I$. In a central Tannakian
condensation the tensor quotient sends $g$ to the unit, which would instead
force $S_g=I$. Thus this condensation route is impossible.

An exact search through all 128 unit-containing subsets shows that the
only proper nontrivial fusion subring of indices 1055,1768,1795 is the
pointed $C_2$ subring. Indices 3258,3413,3472 have no proper nontrivial
fusion subring. The irrational total dimensions therefore rule out every
nontrivial faithful finite-group grading: its neutral component would
have integral dimension one or two. A nontrivial equivariantization
contains $\operatorname{Rep}(G)$; the same subring audit excludes it for
the latter three rings and leaves only $G=C_2$ for the former three.

For 1768 and 1795 the sign invertible acts freely on all eight simples.
The orbit/projective-stabilizer description of equivariantization simples
then forces a hypothetical $C_2$-de-equivariantization to have rank four:
a free base orbit would give a sign-fixed simple, while a fixed base simple
gives two equivariant lifts. Its dimension would be respectively
$222+36\sqrt{37}$ or $270+120\sqrt5$, both greater than $45$.
The unrestricted rank-four companion theorem bounds that dimension by
$24+12\sqrt3<45$. This explicitly cited theorem dependency rules out both
routes; its source digest and the exact quotient tensors are pinned in
the evidence. Together with the $S_g$ argument above, all six residuals
are excluded from nontrivial finite-group equivariantization constructions.
No central lift is inferred from commutativity of the fusion ring alone.

There is also an unbounded arithmetic sieve for generalized-near-group
construction towers. If a base has $G$ invertibles, stabilizer order $s$,
noninvertible orbit size $\ell=G/s$, common irrational dimension $d$, and
total noninvertible multiplicity $K$, then
\[
 d^2=s+Kd,\qquad D_0=\ell(2s+Kd),\qquad
 \frac{\operatorname{N}(D_0)}{\operatorname{Tr}(D_0)}=G.
\]
Pointed products, finite-group extensions, and equivariantizations multiply
dimension by a positive integer $k$, while Morita equivalence preserves it.
Consequently $G_D=\operatorname{N}(D)/\operatorname{Tr}(D)$ must be integral,
$k\mid G_D$, $s\mid G_D/k$, and
$K^2=\operatorname{Tr}(D)/(k\ell)-4s$. These divisor and square tests
exhaust the stated family without a parameter cutoff. Division by
de-equivariantization is outside its scope.

For 1795,3258,3413, the values $G_D=10/3,45/7,45/7$ exclude the whole
stated family. For 1055 and 1768, $G_D=6$ and all matches have $s=1,K=12$;
the tuples $(k,G,\ell)$ are
\[
 (1,6,6),\quad (2,3,3),\quad (3,2,2),\quad (6,1,1).
\]
The last is excluded by the rank-two classification and the third by the
rank-four bound, its base dimension being $148+24\sqrt{37}>45$.
For 3472, $G_D=7$, $s=1,K=14$, and the matches are $(1,7,7),(7,1,1)$;
the latter is again rank-two impossible. The remaining matches are
dimension constraints only, without an asserted base category or Morita
equivalence.

\subsection{Exact character-field coverage}

For the seven indices below, the source element
$H=\sum_{j=0}^7(j+1)^2X_j$ has squarefree degree-eight characteristic
polynomial and generates the rational algebra. Exact interpolation writes
every basis element as a rational polynomial in $H$. Factoring its
polynomial over the quadratic FP field therefore proves the following
full-character coverage, independently of numerical root approximations.

\begin{center}
\begin{tabular}{r|r|r|l}
Index & FP radicand & Characters in FP field & Other factor degrees\\\hline
65 & 5 & 4 & $2,2$\\
1055 & 37 & 6 & $2$\\
1768 & 37 & 8 & none\\
1795 & 5 & 8 & none\\
3258 & 5 & 8 & none\\
3413 & 5 & 5 & $3$\\
3472 & 2 & 2 & $3,3$
\end{tabular}
\end{center}

Thus a modular-eigenvalue calculation confined to the real FP field has
complete character coverage for 1768,1795,3258. The other four need larger
splitting fields. The following exact adapters now supply all four of
them, with every source character polynomial checked to split completely:
\[
\begin{array}{c|c|l}
\text{index}&[L:\mathbb Q]&\text{defining polynomial of }u\\\hline
65&4&u^4-8u^2+36\\
1055&4&u^4-68u^2+1600\\
3413&6&u^6-21u^4+2u^3+84u^2+24u-19\\
3472&6&u^6+2u^5-9u^4-14u^3+10u^2+8u+1
\end{array}
\]
For 65 and 1055, the fields are respectively
$\mathbb Q(\sqrt5,i)$ and $\mathbb Q(\sqrt{37},\sqrt{-3})$, with
$u$ the sum of the displayed radicals. Complex conjugation fixes the
real radical and negates the imaginary one; its multiplicativity and
involution identities are checked exactly. For 3413 and 3472, take
$u=\sqrt q+\theta$, with respectively
$(q,p_\theta)=(5,t^3-3t+1)$ and $(2,t^3+t^2-2t-1)$.

For $\lambda\in L$, algebraic integrality is equivalent to the rational
matrix of multiplication by $\lambda$ having monic integral characteristic
polynomial: necessity follows from the minimal polynomial, and sufficiency
because that characteristic polynomial annihilates $\lambda$.
This tests the full ring of integers without assuming that the primitive
power basis is integral. The adapters check all eight character rows,
their products, dualities, orthogonality, the FP embedding, and direct
multiplication-operator controls. All twelve distinct restriction columns
of the known center of Fibonacci $\boxtimes\operatorname{Vec}(C_2\times C_2)$
are retained in both overfield adapters. These are complete arithmetic
applicability results; categorical obstruction claims retain the
hypotheses of the associated center theorem.

% Follow this fragment by connected_algebra_section.tex, which contains
% the general connected-algebra proof and the complete enumeration table.

\subsection{A complete underlying-object filter for connected algebras}

The following necessary conditions require neither pivotality nor unitarity.
Let $A$ be a connected separable algebra, $m=A$ its simple regular right
module, and $a_i=[A:X_i]$. Free-module adjunction and rigidity identify
$a_i$ with $\dim\operatorname{Hom}(X_i\otimes m,m)$ and then with
$\dim\operatorname{Hom}(m,X_i^*\otimes m)$. Semisimplicity permits reversing
the two arguments without changing this dimension. Hence $a_i=a_{i^*}$.
The positive Frobenius--Perron eigenvector of the module action gives
$a_i\leq\operatorname{FPdim}(X_i)$. If an invertible $g$ occurs in $A$,
the nonzero map $g\otimes m\longrightarrow m$ is an isomorphism, so
$gA=A$ at the level of underlying objects.

Write $B$ for the integral multiplicity matrix of free induction.
Adjunction gives $B^{\mathsf T}B=N_A$, where $N_A$ denotes right
multiplication by $A$ (the source rings below are commutative).
The row corresponding to $m$ is $a$, so
\[
 N_A-aa^{\mathsf T}\quad\text{is entrywise nonnegative and positive semidefinite.}
\]
Thus all underlying objects of connected separable algebras occur in the
finite domain $a_0=1$, $a_i=a_{i^*}$,
$0\leq a_i\leq\lfloor\operatorname{FPdim}(X_i)\rfloor$.
Exhausting this entire domain, imposing invertible stabilization, and
applying the two Gram conditions gives the following complete filter.

\begin{center}
\begin{tabular}{r|r|r|r}
Source index & Bounded self-dual vectors & Entrywise survivors & PSD survivors\\\hline
1055 & 794976 & 35913 & 85\\
1768 & 9539712 & 249139 & 115\\
1795 & 12984608 & 416274 & 87\\
3258 & 580608 & 32098 & 99\\
3413 & 580608 & 35404 & 60\\
3472 & 168591360 & 4841427 & 70
\end{tabular}
\end{center}

The evidence stores every surviving coefficient vector. Integer
fraction-free elimination has a proved arithmetic bound; independent
rational elimination checks every survivor. A complete independent
392-vector control domain for
$\operatorname{Rep}(C_{13}\rtimes C_6)$ reproduces all 30 survivors,
including the regular function algebra with zero residual Gram matrix.
For source index 3472, only $1+X_3$ among all $1+X_i$, $i>0$, survives.
The unit algebra and all internal-end candidates from the rank-two
based-module search survive. A surviving object need not admit an algebra
structure: integral Gram factorization and the module/algebra
associativity equations remain further requirements.

\subsection{Real $STS$ constraints from a quadratic CM character field}
Suppose the complete character field is a CM extension
$L=E(\sqrt{-r})$ of degree two over the Frobenius--Perron field
$E=\mathbb Q(\sqrt q)$. Write $\pi(x)=(x+\overline{x})/2$.
For distinguished central simples the modular identity is
\[
 \sum_Z\theta_Z\chi_A(v_Z)\chi_B(v_Z)
   =f_B\chi_A(v_B).
\]
Apply the normalized trace to $E$ from a common field containing $L$
and the twists. Distinguished twists are one. Complete high dual pairs
have equal twists and conjugate character products, so their summed
coefficient lies in $E$ and factors through the existing FP-field
relative twist trace. Consequently the low choice $(v,G)$ has coefficient
\[
 \pi\!\left(\chi_A(v)\chi_B(v)
       -\delta_{G,B}f_B\chi_A(v)\right),
\]
with right-hand side $\pi(\delta_{B,\mathrm{FP}}f_B-1)$.
Thus these additional real $STS$ rows require no enlarged phase field.
The separate full-character-field $S$-symmetry equations are retained.

The isolated implementation checks conjugation invariance of every
high-orbit coefficient and reconstructs every coefficient in the exact
FP basis. Its independent
$\mathrm{Fib}\boxtimes\mathrm{Vec}_{C_4}$ control checks an exact feasible
primal, all modular columns and all factored-oracle scores against
direct arithmetic. This strengthening applies to $\sha{R8_3_65}$, whose
field is $\mathbb Q(\sqrt5,i)$. It does not apply to $\sha{R8_3_3472}$:
there the complete real character field is strictly larger than the FP
field, and an $STS$ calculation would require a separate extension of
the relative-trace domain.

For $\sha{R8_3_65}$, the complete height-four relaxation has 283 rows and
6521 columns. It admits an exactly verified nonnegative rational
solution. Thus this additional modular constraint is exhausted at the
stated height and does not exclude the ring; the feasible relaxation
does not construct a fusion category. Both the positive-control receipt
and the target primal replay pass with Python assertions disabled.

\subsection{A complete noncommutative Galois--pseudo-unitary reduction}
\label{subsec:frontier-nc-galois-pu}

\begin{proposition}
Every complex fusion categorification of $\sha{R8_3_721}$ or $\sha{R8_3_1154}$
admits a spherical structure and has a pseudo-unitary Galois conjugate.
\end{proposition}
\begin{proof}
For each ring, the complexified fusion algebra is
$\mathbb C^4\oplus M_2(\mathbb C)$. The common kernel of all commutators
in the regular representation consists exactly of the scalar blocks.
The four characters obtained there exhaust all possible dimension
characters; a matrix block has no one-dimensional representation.
Its formal codegree is nevertheless included: if its regular Casimir
eigenvalue is $\lambda$, with multiplicity four, its formal codegree is
$\lambda/2$. Exact multiplication identities and the complete Casimir
factorization verify this accounting. For each original ring, every scalar-character orbit
except the FP orbit is excluded by a zero dimension or formal-codegree
norm divisibility.

For an arbitrary categorification take its canonical sphericalization~\cite[Section~5.1]{OstrikPiv3}
and choose dual-compatible lifts. The signed quotient by the sign object
$\varepsilon=-1$ satisfies the full parity intervals
$|T_{ij}^k|\leq N_{ij}^k$ and $T_{ij}^k\equiv N_{ij}^k\pmod2$,
cyclic trace symmetry, dual reversal, and every ordered associativity
equation. We impose neither commutativity nor a modulo-four indicator
normalization. Complete enumeration gives $64$ signed tensors for
$\sha{R8_3_721}$, all gauge-equivalent to the original tensor, and $288$
for $\sha{R8_3_1154}$. The latter comprise one $64$-member original gauge
class and five other classes of sizes $64,64,32,32,32$.
All scalar-character factors of the five other classes are excluded by
zero dimensions or exact norm divisibility. Both sign sectors' formal
codegrees, including matrix blocks, are used: sphericalization has
dimension $2\delta$ and formal codegrees $2f$, so central induction
requires $\delta/f$ to be integral~\cite[Theorem~2.13 and Corollary~2.14]{OstrikPiv3}.

Thus every possible signed quotient is gauge-original. Switching the
lifts by its dual-compatible gauge makes the signed tensor equal to
$N$. The negative lift multiplicities are then zero, so the chosen
lifts form a spherical fusion subcategory whose forgetful functor to
the original category is a tensor equivalence. Its dimension character
lies in the FP Galois orbit. A number-field model~\cite[Theorem~2.30 and the following remark]{ENO} and Galois conjugation
therefore give a spherical category with positive FP dimensions, hence
a pseudo-unitary category with the same fusion rules.
\end{proof}

The source-bound replay is
\path{central_integrality/general_galois_pu_gate.py} in the frontier
evidence directory; its primary calculation incorporates the independent
unrestricted signed enumeration. The earlier normalized enumeration is
retained only as a cross-check. Consequently the central-support, partial
modular, integral Gram, joint-rotation, and low Verlinde constraints for
these two rings are necessary for general complex existence after Galois
transport. Passing these constraints supplies neither a model nor
unitarity; both rings remain unresolved by this proposition.

\subsection{Spherical gates for all seven noncommutative frontiers}
\label{subsec:frontier-nc-quartic-and-tests}

\begin{proposition}
Every complex fusion categorification of any of
$\sha{R8_3_439}$, $\sha{R8_3_455}$, $\sha{R8_3_567}$, $\sha{R8_3_721}$,
$\sha{R8_3_928}$, $\sha{R8_3_941}$, and $\sha{R8_3_1154}$
admits a spherical structure and has a pseudo-unitary Galois conjugate.
\end{proposition}
\begin{proof}
The quadratic cases are proved in
Section~\ref{subsec:frontier-nc-galois-pu}.
The same unrestricted signed construction gives the following complete
domains in the two quartic cases:
\[
\begin{array}{c|r|r|r}
\text{ring}&\text{ordered signed tensors}&\text{gauge-original}&
 \text{nongauge excluded}\\\hline
R8\_3\_928&32&16&16\\
R8\_3\_941&128&64&64
\end{array}
\]
No commutativity or modulo-four indicator restriction is imposed.
Each nongauge part is one dual-compatible sign-gauge class. Its
spherical dimension has polynomial
$t^4-320t^3+18688t^2-139264t+278528$, whereas the ordinary scalar
formal codegrees have polynomial
$t^4-576t^3+11520t^2-73728t+147456$.
Induction in the sphericalization requires their quotient to be an
algebraic integer. In a common number field its norm is a positive
integer power of $278528/147456=17/9$, a contradiction.

The scalar algebra of each ordinary ring is the irreducible quartic
FP orbit, and the full complex algebra is
$\mathbb C^4\oplus M_2(\mathbb C)$.
An independent check constructs all five central idempotents and verifies
regular ranks $(1,1,1,1,4)$. For the nongauge signed tensors, rational
regular projectors verify a separable quartic scalar algebra and a
four-dimensional matrix algebra with one-dimensional center.
Its Casimir eigenvalue is $8$, hence its formal codegree is $4$.
Thus no dimension character or matrix formal codegree is omitted.
All surviving signed tensors are gauge-original; the spherical-section
and Galois-transport argument of the preceding proposition applies.

For $439,455,567$, the unrestricted signed domains have respectively
$256,32,64$ tensors, of which $64,32,32$ are gauge-original. Rational
regular projectors again verify every scalar and matrix component,
including the matrix formal codegrees. The last nongauge class of both
$439$ and $567$ has a rational scalar character of dimension $\delta=24$.
In a sphericalization with dimension field $K$, every central simple
dimension belongs to $K$, since its forgetful image is a sum of simples.
For an ordinary-sector formal codegree $f$, induced-unit decomposition
therefore gives $\dim(A_E)=(2\delta)/(2f)=\delta/f\in K$; see
\cite[Theorem~2.13 and Corollary~2.15]{OstrikPiv3}. Hence $f\in K$.
Here $K=\mathbb Q$, whereas an ordinary formal codegree has irreducible
polynomial $t^2-96t+576$. This excludes the last nongauge class without
assuming pseudo-unitarity. All other nongauge classes are excluded by
zero, duality or norm conditions. The ordinary scalar-character sieve
leaves precisely the FP orbit, completing the same spherical-section
and Galois-transport proof.
\end{proof}

The canonical quartic replay is
\path{noncommutative_followup/verify_scope_gates.py}.
The additional three-case replay is
\path{noncommutative_followup/additional_scope_gates.py}.
Together the seven gates make the following pseudo-unitary calculations
necessary for general complex existence after Galois transport; they
make no assertion of unitarity or existence.

Write $B=MM^{\mathsf T}$ for the induction Gram matrix and the central
restriction matrix. Every restriction column is orthogonal to
$\ker B$. Imposing this exact central-image condition before enumeration
leaves respectively $13,1,1,59$ weighted distinguished frames for
$721,928,941,1154$; ordinary Gram and moment bounds leave $1,1,1,4$.
The two-dimensional representation contributes one distinguished column
with unit entry $2$, whose outer square is counted once.

For an irreducible trace character $\chi_E$, put $m_E=\dim E$.
The distinguished central simple $A_E$ has twist $1$, and its complete
unnormalized $S$ row is given by \cite[Equation~(39)]{DeshpandeCenter}:
\[
 S_{A_E,Z}=\frac{D}{f_E m_E}\chi_E(FZ),\qquad D=\operatorname{FPdim}(\mathcal C).
\]
The matrix-block character is included. Symmetry and
$STS=DT^{-1}ST^{-1}$ impose
\[
 \sum_Z\chi_E(FZ)\chi_F(FZ)\theta_Z
 =f_Fm_F\chi_E(FA_F).
\]
These supplement the linear indicator equations and the joint
divisor-power rotation constraints. They remove one of two symmetric
assignments for $721$ and one of the last two frames for $1154$.
The surviving $928$ and $941$ branches have exact rational witnesses.
The central eigenvalues $S_{A_E,Z}/\dim Z$ are also tested for algebraic
integrality.

For $721$, complete integral Gram enumeration has four patterns.
Each survives joint rotation through height $12$ and the combined
height-$6$ rotation and $STS$ system with its multiplicities fixed.
For selfdual distinguished simples, the ordinary Verlinde formula gives
\[
 N_{A_E,A_F}^{A_G}
 =D^{-2}\sum_Z
 \frac{S_{A_E,Z}S_{A_F,Z}S_{A_G,Z}}{\dim Z}.
\]
All twenty coefficients are nonnegative integers for every pattern;
the corresponding forgetful tensor equations remain exactly feasible.
For $1154$, the complete integral Gram domain has $14211$ patterns.
Cubic Verlinde integrality leaves $1238$; all pass the ten exact rational
forgetful tensor cones. The subsequent fixed-pattern rotation calculation
is given separately below.

The exact pointed-bimodule stabilizer test also admits both $C_2$
associator classes wherever applicable. The forced pseudo-unitary
invertible indicators are $-1,+1,-1$ for $928,941,1154$, respectively.
These checks do not construct the required module coherences or
associators. Four higher-height numerical attempts remain explicitly
undecided and are not counted as completed obstructions. Exact feasible
relaxations, including those above, are not categorical models.

\subsection{An all-level quantum-group product and Morita-family sieve}
\label{subsec:frontier-all-level-su2}

\begin{lemma}
Let $a_1,\ldots,a_r$ be totally positive algebraic numbers, each chosen
as a largest conjugate. Then
$\mathbb Q(\prod_i a_i)$ is the compositum of the fields $\mathbb Q(a_i)$.
\end{lemma}
\begin{proof}
In a common Galois closure, an automorphism fixing the product has
$0<\sigma(a_i)\le a_i$ for every $i$. Equality of products forces equality
in each factor. Thus the product stabilizer is the intersection of the
individual stabilizers, and the assertion follows from Galois theory.
\end{proof}

For $n=k+2\ge3$, the adjoint $\mathrm{SU}(2)_k$ category has global
dimension
\[
 A_n=\frac{n}{2-2\cos(2\pi/n)}.
\]
The full category has dimension $2A_n$; see
\cite[p.~8]{CoquereauxQuantumMcKay}
and the equal-component theorem \cite[Proposition~8.20]{ENO}.
All conjugates of $A_n$ are positive, $A_n$ is largest, and
$[\mathbb Q(A_n):\mathbb Q]=\varphi(n)/2$.
Consequently a proposed identity $D=m\prod_n A_n^{e_n}$, with
$m\in\mathbb Z_{>0}$, forces every factor field into $\mathbb Q(D)$.
The verifier uses the possibly larger FP field, a safe enlargement of
the search domain, and verifies its normality exactly.

If this field has degree $d$, then $\varphi(n)\mid2d$.
Every prime divisor $p$ of $n$ satisfies $p-1\mid2d$, and its exponent
is bounded by $\varphi(p^a)\le2d$. These conditions give a complete
finite domain with no level cutoff. Exact real cyclotomic resultants
and factorization check field containment. Rational factors
$A_3=1$, $A_4=2$, $A_6=6$ are absorbed into $m$; all remaining factors
exceed one. The inequalities $A_n^{e_n}\le D$ bound all exponents
exactly, after which every product and integer quotient is checked.

The replay \path{noncommutative_followup/su2_all_frontier.py}
gives $30$ family exclusions among the $34$ original frontier rings.
The retained dimension matches are precisely
\[
\begin{array}{c|c}
R8\_3\_439,\ R8\_3\_567&D=2A_{12}\\
R8\_3\_731&D=660\\
R8\_3\_3344&D=360.
\end{array}
\]
The two integral cases use the empty product. None of these matches
is a constructed category. The center-invertible obstruction in
Section~\ref{subsec:frontier-known-family-closure} subsequently excludes
the two nonintegral matches throughout the stated categorical operation
family, giving $32$ family exclusions in total.

The excluded family includes arbitrary finite Deligne products of the
full and adjoint $\mathrm{SU}(2)$ categories at every positive level,
followed by any finite sequence of finite-group extensions,
equivariantizations, products with integer-FP-dimension categories, and
Morita equivalences. The first three operations multiply global
dimension by a positive integer; Morita equivalence preserves it
\cite[Corollary~8.14]{ENO}. The statement excludes this construction
family only. Other Lie types, arbitrary algebra condensation,
dimension-decreasing de-equivariantization, and additional nonintegral
product factors are outside its scope.

The generalized near-group dimension argument also extends uniformly
to the $34$ original frontier rings. Write $G=\ell s$ for the invertible
group order, $\ell$ for the noninvertible orbit size, and $s$ for the
stabilizer order. If $d^2=s+Kd$, then an irrational positive integer
dimension multiple $D=k\ell(s+d^2)$ is quadratic and satisfies
\[
 \frac{\operatorname{Norm}(D)}{\operatorname{Trace}(D)}=kG,
 \qquad K^2=\frac{\operatorname{Trace}(D)}{k\ell}-4s.
\]
Thus $k$ divides the forced integer $G_D=\operatorname{Norm}(D)/
\operatorname{Trace}(D)$, and $s$ divides $G_D/k$. Exhausting these
divisors gives every possible parameter tuple, with no rank cutoff.
The exact replay agrees with all six earlier quadratic computations.
It yields $20$ dimension-family exclusions, $12$ complete finite
parameter lists, and two integral targets outside this quadratic test.
The parameter lists are not models. The family here has a single
generalized near-group source and a positive integer multiplier;
products of distinct nonintegral sources and dimension-decreasing
operations are not included. The canonical table is
\path{noncommutative_followup/generalized_near_group_all_frontier.py}.

\subsection{Earlier central tests for three further noncommutative rings}
\label{subsec:frontier-additional-nc-central}

The same central-image and full-character $S$ symmetry computation has
$5,1,5$ initial frames for $439,455,567$ and leaves one frame in each
case, with two symmetric assignments. For $567$ the scalar characters
do not all lie in the real FP field; the computation adjoins $i$ and
verifies complete character orthogonality in the resulting exact
quartic field. No character is discarded for being nonreal.

On the unique frame for each of $439$ and $567$, all eight second-indicator
branches are tested against joint tensor rotations through height six.
Fourteen exact dual certificates exclude seven branches per ring, and
two exact primal witnesses certify the remaining relaxations. In the
original basis, the surviving tuples for objects $1,\ldots,7$ are
\[
 \nu_2^{439}=(0,0,1,1,1,1,1),\qquad
 \nu_2^{567}=(0,0,0,0,-1,-1,-1).
\]
In particular, the three odd $S_3$ invertibles have indicator $+1$ in
$439$ and $-1$ in $567$, in the pseudo-unitary Galois representative.
The canonical replays are:
\begin{quote}\small
\path{noncommutative_followup/additional_central_constraints.py}\par
\path{noncommutative_followup/additional_joint_rotation.py}.
\end{quote}
The surviving frames and witnesses remain necessary conditions, not
categorical constructions.

\begin{proposition}[A two-prime norm obstruction for integral subcategories]
\label{prop:frontier-integral-center}
Let $\mathcal B$ be a fusion category with no nontrivial invertible simple
objects. Suppose that its integral-dimensional simple objects span a fusion
subcategory $\mathcal E$, and that the absolute field norm of
$\operatorname{FPdim}(\mathcal B)$ has at most two distinct prime divisors.
Then $\mathcal E=\mathrm{Vec}$.
\end{proposition}
\begin{proof}
Write $N=\operatorname{FPdim}(\mathcal E)\in\mathbb Z_{>0}$ and
$K=\mathbb Q(\operatorname{FPdim}(\mathcal B))$.
Subcategory divisibility \cite[Proposition~8.15]{ENO} implies that
$\operatorname{FPdim}(\mathcal B)/N$ is an algebraic integer.
Taking its norm shows that every prime divisor of $N$ divides the norm of
$\operatorname{FPdim}(\mathcal B)$. Thus $N=p^a q^b$ for at most two primes.
The category $\mathcal E$ is solvable by \cite[Theorem~1.6]{ENOSolvable}.
If nontrivial, it has a nontrivial invertible simple by
\cite[Proposition~4.5(iv)]{ENOSolvable}, contradicting the hypothesis on
$\mathcal B$.
\end{proof}

For any hypothetical category $\mathcal C$ with ring $\sha{R8_3_439}$ or
$\sha{R8_3_567}$, the already verified center-invertible obstruction gives
$\operatorname{Inv}Z(\mathcal C)=\{\mathbf1\}$.
Its simple dimensions are six copies of $1$ and two copies of
$d=3+2\sqrt3$, so $D=\operatorname{FPdim}(\mathcal C)=48+24\sqrt3$.
Every central simple restricts to an object of dimension $a+bd$ with
$a,b\in\mathbb Z_{\geq0}$. It has integral dimension precisely when $b=0$.
These objects form a fusion subcategory: their forgetful images lie in the
pointed source subcategory, which is closed under tensor products and duals.
Since $\operatorname{FPdim}Z(\mathcal C)=D^2$
\cite[Proposition~8.12]{ENO}, its integral subcategory dimension $N$ satisfies
\[
 \operatorname{Norm}_{\mathbb Q(\sqrt3)/\mathbb Q}(D^2/N)
 =\frac{576^2}{N^2}\in\mathbb Z,
 \qquad N\mid576=2^6 3^2.
\]
Proposition~\ref{prop:frontier-integral-center} therefore shows that the tensor
unit is the only integral-dimensional central simple. This conclusion requires
no pivotal or pseudo-unitary hypothesis.

The same center condition holds for the $R_{5,3,40}$ Morita dual forced by
the full-$S_3$ algebra in a hypothetical $\sha{R8_3_439}$ category. Its source
dimensions are $1,1,2,d,d$, and the integral source objects form the
$\operatorname{Rep}(S_3)$ based subcategory. The same proof removes every
nonunit central restriction column supported on $1,s,V$. This assertion is
about the specified Morita dual; it does not impose trivial center invertibles
on arbitrary categorifications of $R_{5,3,40}$.

\subsection{Complete fixed-Gram rotation tests for $\sha{R8_3_1154}$}
\label{subsec:frontier-fixed-gram-1154}

For $\sha{R8_3_1154}$, the complete central-support and low-frame calculation
leaves frame $31$, the second-indicator tuple $(0,0,-1,1,1,1,1)$,
and sixteen high restriction types. Exact affine bounds enumerate all
$14{,}211$ nonnegative integral Gram completions; every completion is
checked against the original $28$ Gram equations. Of these, exactly
$1{,}238$ have nonnegative integral distinguished low Verlinde
coefficients, and all of these also satisfy the quadratic spectral
multiplicity conditions and the ten fixed low-product rational cones.

We next fix each of these $1{,}238$ integral multiplicity vectors in the
joint rotation system through order six, together with all distinguished
character $STS$ equations. Four exact separating rays exclude eight
vectors. The remaining $1{,}230$ vectors have exact nonnegative rational
solutions. A source-bound receipt covers every vector in $434$
certificate families: each separating ray is checked on the complete
column domain, and every primal family is checked in the original
matrix equation at each vector assigned to it. No pattern is left to
a numerical feasibility decision. The canonical read-only replay is
\path{central_integrality/fixed_gram_rotation_1154.py}.

The Galois--pseudo-unitary reduction makes this a necessary-condition
calculation for general complex existence after Galois transport.
The finite bound is rotation order six, not a bound on the categorical
exponent. The surviving rational solutions do not supply central fusion
rules or coherent associators, and $\sha{R8_3_1154}$ remains open.

\subsection{Integer central restriction and cubic constraints}
\begin{proposition}[Integer central necessary system]
Let $\mathcal C$ be pseudo-unitary with commutative Grothendieck ring,
$D=\operatorname{FPdim}(\mathcal C)$, characters $\chi_E$, and formal
codegrees $f_E$. Put $\beta_E=D/f_E$, and let $A_E$ denote the distinguished
simple of $Z(\mathcal C)$ belonging to $E$. A complete finite restriction
column domain gives a finite system of integer equations with bounded
nonnegative integer variables that includes restriction Gram identities,
distinguished $S$-symmetry, duality, and
\[
 N_{A_E,A_F}^{A_G}
 =\sum_Z\frac{\beta_E\beta_F\beta_{G^*}
 \chi_E(FZ)\chi_F(FZ)\chi_{G^*}(FZ)}{D^2\operatorname{FPdim}(Z)}.
\]
Every categorification supplies an integer solution of this system.
\end{proposition}
\begin{proof}
By \cite[Lemma~4.3 and Corollary~4.5]{DeshpandeCenter},
$S_{A_E,Z}=\beta_E\chi_E(FZ)$ in the unnormalized modular convention.
The displayed identity is the Verlinde formula for the center of dimension
$D^2$. Introduce a nonnegative integer selection variable for each allowed
restriction of each distinguished simple, with sum one for each character.
The tensor unit is fixed separately; its character is located by equality
with the positive dimension character. High variables count actual central
simples per complete restriction-duality orbit. Gram, modular symmetry and
low-duality equations are linear in these variables. Expanding the displayed
formula in the full character field makes every nonconstant coordinate zero
and its constant coordinate a nonnegative integer auxiliary variable.
Diagonal Gram bounds and the sum of positive constant contributions give
finite upper bounds. Clearing denominators produces the stated system.
No integrality is imposed on relative twist-trace weights.
\end{proof}
The eigenvalues $\beta_E\chi_E(FZ)/\operatorname{FPdim}(Z)$ of the integer
fusion matrix of $A_E$ also have equal multiplicities throughout each full
Galois orbit. These yield further exact linear equations. Integer interval
and gcd propagation are necessary tests. A covector annihilating all columns
modulo an integer $m>1$, but not the right-hand side, certifies nonmembership
in the column lattice; rational annihilators are checked in the same manner.
Signed integer lattice membership alone is weaker than a nonnegative
bounded solution. Numerical optimization only proposes vectors, which are
checked against every exact original equation and bound. Any negative
pseudo-unitary result still requires a complete signed and dimension-character
argument before it is promoted to a general-complex exclusion.
The actual center of $\mathrm{Fib}\boxtimes\mathrm{Vec}_{C_4}$ provides a
64-simple positive integer control with complex character values.
The dimension identity further gives the exact bound
\[
 N_{A_E,A_F}^{A_G}\leq
 \left\lfloor\frac{\beta_E\beta_F}{\beta_G}\right\rfloor.
\]
If a candidate eigenvalue has a Galois conjugate absent from the complete
remaining domain, its nonnegative multiplicity is zero; this rule may be
iterated, retaining the fixed constituents. Row content division, exact
proportional-row identification, and zero-row deletion preserve the integer
solution set. Their signed row scales also lift any compact-system
annihilator back to a covector checked on every original column.

\subsection{Joint Galois spectra and exact fusion bounds}

The integer central reconstruction can be strengthened without enlarging
the twist field.  Let $A_E$ be the distinguished simple objects in
$I(\mathbf1)$, let $\beta_E=\operatorname{FPdim}(A_E)$, and suppose that
a normal number field $L$ contains the complete character table and the
Frobenius--Perron dimensions.  Write $m_E=\dim(E)$ for the degree of the
irreducible representation affording the full matrix character $\chi_E$.
For a central simple with restriction vector $v$, put
\[
 \Lambda(v)=\left(
 \frac{\beta_E\chi_E(v)}{m_E\operatorname{FPdim}(v)}
 \right)_{E\ne\mathrm{FP}}.
\]
In a pseudo-unitary categorification, the multiset of these complete
tuples, counted with central-simple multiplicities, is invariant under
$\operatorname{Gal}(L/\mathbb Q)$.  In the commutative implementations
here, every $m_E=1$; the degree factor is required for full characters
of noncommutative source rings.

Indeed, the distinguished fusion matrices are commuting integral
matrices in the braided center and are simultaneously diagonalized by
its modular matrix.  Choose an integer linear form separating the
finite union of the candidate tuples and all their Galois conjugates.
Such a form exists because only finitely many proper rational
hyperplanes are forbidden.  The corresponding integer linear
combination of the fusion matrices has an integral characteristic
polynomial.  Equality of multiplicities of its conjugate eigenvalues,
together with separation, proves the assertion for complete tuples.
The distinguished labels remain fixed; the same field automorphism
acts on every tuple coordinate.  No braiding of the source category is
assumed.

Consequently, complete-tuple multiplicities give additive integer
equations in the low choices and high duality-orbit counts.  Fixed
constituents contribute to the right-hand side.  A tuple whose
conjugate is absent from the remaining candidate support forces every
nonnegative variable containing that tuple to vanish; this deletion
may be iterated.  Missing conjugates of a fixed constituent instead
give an obstruction.  These joint equations strengthen the separate
Galois multiplicity equations for individual distinguished matrices.

The nonnegative integer variables recording distinguished fusion
coefficients also satisfy the exact bounds
\[
 0\le N_{A_E,A_F}^{A_G}
 \le \left\lfloor\frac{\beta_E\beta_F}{\beta_G}\right\rfloor.
\]
The floor is evaluated in the positive Frobenius--Perron embedding by
exact sign comparisons.  The joint-spectrum and bound adapters retain
the actual $64$-simple center of $\mathrm{Fib}\boxtimes\mathrm{Vec}(C_4)$,
including its nonreal characters, in normal and assertion-disabled
replays.  These are necessary reconstruction conditions: a surviving
integer or rational solution does not supply an associator, while a
numerical infeasibility report alone is never an exclusion certificate.

\subsection{A pointed integral-center criterion}
\begin{proposition}
Let $\mathcal C$ be pseudo-unitary. Suppose its simple dimensions are
$a_i+b_i\sqrt q$, with $b_i\geq0$, and $b_i=0$ precisely for invertible
simples. If the tensor unit is the only integral-dimensional simple
constituent of $I(\mathbf1)$, then the integral-dimensional subcategory
of $Z(\mathcal C)$ is pointed, and the forgetful functor is injective
on its invertible objects.
\end{proposition}
\begin{proof}
An integral-dimensional central simple restricts to the pointed source
subcategory, since the irrational coefficients of source dimensions are
nonnegative. By adjunction, such a simple can contain the source unit
after forgetting only if it is itself the unit. For an integral central
simple $Z$, every constituent of $ZZ^*$ is integral: its restriction
belongs to the pointed source subcategory. Hence
\[
 1=[\mathbf1:ZZ^*]=[\mathbf1:F(ZZ^*)]
   =\sum_{g\in G(\mathcal C)}[g:F(Z)]^2.
\]
Thus $F(Z)$ is a single invertible and $Z$ is invertible. Two such
invertibles with the same restriction have quotient restricting to the
unit, so their quotient is the unit.
\end{proof}
For commutative source rings, the hypothesis on $I(\mathbf1)$ is checked
on the complete distinguished dimension list $D/f_E$. The criterion
removes every integral high restriction vector other than a single
source invertible. Its exact quadratic hypotheses and every deleted
column are recorded for $\sha{R8_3_1768}$ and $\sha{R8_3_1795}$.

\subsection{Source-unit identities for distinguished products}
Let the source ring be commutative, and write $v_E=[F(A_E)]$ for a
distinguished central simple. Its multiplicity in $I(\mathbf1)$ is
$m_E=1$, so $(v_E)_0=1$; every other central simple has source-unit
restriction multiplicity zero. Write $E^*$ for the dual character label,
and let $v^*$ denote source duality on a restriction vector. Taking the
source-unit coefficient of $F(A_EA_F)$ gives the exact identity
\[
 \delta_{F,E^*}+\sum_{G\ne\mathrm{FP}}N_{A_E,A_F}^{A_G}
 =\langle v_E,v_F^*\rangle.
\]
In particular, for each of the seven nonunit distinguished labels,
\[
 1+\sum_{G\ne\mathrm{FP}}N_{A_E,A_{E^*}}^{A_G}=\|v_E\|^2.
\]
The right side is linear in the atomic one-hot selection variables:
$\|v_E\|^2=\sum_v x_{E,v}\|v\|^2$. For the other unordered pairs,
\[
 2\left(\delta_{F,E^*}+\sum_{G\ne\mathrm{FP}}N_{A_E,A_F}^{A_G}\right)
 \leq\|v_E\|^2+\|v_F\|^2,
\]
since the difference is $\|v_E-v_F^*\|^2$. These become exact integer
equations by adjoining nonnegative integer slack variables. Their upper
bounds are the sums of the largest allowed squared restriction norms.
For a fixed complete frame the norms are constants. Thus the same adapter
applies to atomic choices and complete frames, including nonselfdual
character labels. These rows can strengthen a fractional atomic
relaxation; they are necessary consequences of an actual category and are
not additional existence assertions. The actual 64-simple center of
$\mathrm{Fib}\boxtimes\mathrm{Vec}_{C_4}$ satisfies every new identity
and bound in exact arithmetic, with both normal and optimized replay.

\paragraph{Compatibility of integer restrictions with phases.}
An integer solution of the preceding system and a rational solution of the
rotation and distinguished-$STS$ cone need not describe the same restriction
multiplicities. We therefore also specialize the phase cone to selected
integer solutions. Fix each selected low column with coefficient one, retain
all relative-phase columns for every high restriction-duality orbit $O$ of
positive multiplicity $n_O$, and add
\[
 \sum_{\tau}x_{O,\tau}=n_O.
\]
All other high columns vanish, while proper-divisor rotation columns remain.
Every categorical center with this restriction pattern satisfies the
specialized system. The compressed trace weights $x_{O,\tau}$ remain
nonnegative rational variables; they are not assumed integral. Exact primals
are also expanded back into the original full phase cone, and exact
separators are checked on every specialized column. Rejection of a selected
integer pattern does not exclude other integer patterns or the source ring.

\paragraph{Complete fixed-frame arithmetic screen for \sha{R8_3_982}.}
The distinguished-$S$ constraints leave exactly four labeled low frames.
For each frame we enumerate every high restriction vector allowed by the
residual Gram form, then impose integer Gram, distinguished cubic,
Galois-balance and exact Frobenius--Perron bounds.  The four signed column
lattices contain their right-hand sides.  Two frames admit explicitly
checked nonnegative integer solutions, each with $25$ positive high
restriction-duality orbits and proposed center rank $56$.  Both solutions
also satisfy every joint-Galois balance equation and all $28$ exact
source-unit product identities among the seven nontrivial distinguished
objects.

These two integer solutions do not survive the phase compatibility test.
For each, the height-$4$ specialization has $366$ equations and $885$
variables; exact separators cover all eight indicator choices.  This
excludes the two displayed restriction patterns, not other integer
solutions.  For the other two frames, the strengthened bounded continuous
systems have exact rational primals, checked against every original
equation and bound.  Their reduced systems have $504$ equations and $361$
variables.  Numerical integer infeasibility proposals are retained only
as search diagnostics.  Thus this complete four-frame arithmetic screen
and the two complete selected-pattern phase checks leave
\sha{R8_3_982} unresolved; they provide neither a category model nor
a general noncategorifiability theorem.

\subsection{Global phase cuts on integer reconstruction data}

The integer reconstruction and the phase equations can be coupled by
exact separating inequalities.  Fix a complete phase operator, with
low columns $C_l$, high columns $C_{o,t}$ indexed by restriction-duality
orbits and complete relative-trace types, and auxiliary rotation
columns.  An actual integer pattern has selected low counts $x_l$ and
high orbit counts $n_o$; its nonnegative phase weights satisfy
$\sum_t w_{o,t}=n_o$.

If an integer covector $y$ pairs nonnegatively with every auxiliary
column, define
\[
 a_l=yC_l,\qquad m_o=\min_t yC_{o,t}.
\]
For each allowed indicator choice $\varepsilon$, every compatible
integer pattern satisfies
\[
 \sum_l a_lx_l+\sum_o m_on_o\le yb_\varepsilon.
\]
This follows by applying $y$ to the phase equation and bounding each
complete high fiber by its minimum times its total mass.  The minima
are checked over every original orbit, including orbits absent from
the pattern that suggested $y$.  A common inequality valid for all
indicator choices uses $\max_\varepsilon yb_\varepsilon$ on the right.
The low columns retain their character labels.

For nonnegative integer variables, a valid inequality $a\cdot x\le b$
also implies
\[
 \sum_j\left\lfloor a_j/d\right\rfloor x_j
 \le\left\lfloor b/d\right\rfloor\qquad(d\in\mathbb Z_{>0}).
\]
This exact weakening, valid also for negative coefficients, permits
smaller coefficients in a numerical proposal.  The original full-domain
covector, every indicator right-hand side, the divisor, and all rounded
coefficients remain independently checkable.  An integer slack converts
the inequality into an equality, with a finite upper bound obtained
from the existing variable box.

The finite application to $\sha{R8_3_65}$ considers all four indicator
choices and allows at most ten such global cuts.  Every proposed integer
pattern is checked against the complete integer equations, every phase
separator is checked against the complete specialized phase domain, and
each added global cut must strictly reject its originating pattern.
A numerical failure to find the next integer point remains inconclusive;
it is not an exclusion certificate.

The recorded $\sha{R8_3_65}$ campaign added one exact uniform cut.  It
rejects the initial integer reconstruction for all four indicator
choices, with exact violation $512$ after integer rounding.  The
strengthened system has $1476$ equations and $524$ nonnegative integer
variables, including source-unit product constraints and the conditional
pure-invertible restriction bounds.  Its finite-box and modulo-$2,3$
checks pass.  The next bounded integer proposal was inconclusive, so
this result excludes the selected reconstruction and establishes a
global necessary inequality; it does not exclude the fusion ring.

For $\sha{R8_3_3472}$, the final joint-spectrum, cubic, Frobenius--Perron-bound,
and source-unit-product system has $3347$ equations and $13242$ variables.
The complete support and finite-box filters and the modulo-$2,3$ tests
pass.  The bounded linear search produced no exact separator, and the
bounded integer proposal produced no checked witness.  These outcomes
are explicitly inconclusive.  No complete Hermite-lattice or bounded
integer enumeration is asserted, and this character field's larger
real trace-field STS problem remains outside the tested phase domain.

\subsection{Integral central multiplicities and forced pointed lifts}
\label{subsec:frontier-quartic-integer}

The two noncommutative rings $\sha{R8_3_928}$ and $\sha{R8_3_941}$ admit a common
four-coordinate description of their residual central Gram domain:
\[
 Q=\begin{pmatrix}8&8&8&12\\8&50&58&73\\8&58&74&89\\12&73&89&117\end{pmatrix}.
\]
Their unique distinguished induction frame has 65 possible high restriction
columns.  All pass the full-character partial-$S$ algebraic-integrality test.
The twenty distinguished cubic Verlinde expressions coincide in
$\mathbb Q(\sqrt2,\sqrt3)$.  An exact rational reduction followed by two
independent bounded integer enumerations yields precisely 343 solutions of
the full Gram and nonnegative integral cubic equations.  The native search
visits 1665 nodes and the independent Python finite-set search 561 nodes;
every resulting vector is checked against the original equations.

All 343 vectors contain exactly one central simple restricting to the
nontrivial source invertible $g$.  Since each source is its own adjoint
subcategory, the forgetful kernel is trivial.  The corresponding central
simple $G$ is therefore invertible of order two, and its lift over $g$ is
unique.  Tensoring the five distinguished objects by $G$ forces five
additional distinct central simples, with canonical high indices
$40,41,42,43,57$.  All other high restriction types that occur are fixed
as restriction vectors by multiplication by $g$.  This proves a statement
about the center, without asserting a tensor-factor decomposition of the
nonbraided source.

For $\sha{R8_3_928}$ the already certified positive-spherical indicator is
$\nu_2(g)=-1$, so $\theta_G=\pm i$.  Monodromy with an order-two invertible
has values $\pm1$.  Balancing consequently prevents $G$ from fixing a
central simple.  Multiplicities of the fixed restriction types must be
even, leaving exactly 15 of the 343 vectors.  Their paired second and
sixth twist moments vanish; the five translated distinguished objects
have second and sixth moments $-1$.  The subsequent finite systems retain
the actual correlation between the two twists in each semion orbit.

For $\sha{R8_3_941}$ the indicator is $+1$, and $\theta_G=\pm1$.
The positive-twist branch would give central $\operatorname{Rep}(C_2)$.
The sign character acts freely on all eight source simples, so
de-equivariantization has four simples.  Its exact quotient tensor is
the historical rank-four ring $F_{6,19}$, and its dimension is
\[
 72+48\sqrt2+40\sqrt3+28\sqrt6>72.
\]
This contradicts the unrestricted rank-four companion bound
$24+12\sqrt3<45$ (Corollary~\path{finite:cor:sharp-bounds}).
Hence $\theta_G=-1$: the lift is a fermion.  The ordinary rank-four bound
has been used only in the Tannakian branch.

Write $q_Z=\widetilde S_{G,Z}/d_Z\in\{\pm1\}$.
Positive-charge simples occur in distinct pairs with opposite twists;
negative-charge simples may be fixed.  For each fixed restriction type
write $2p_v+q_v=m_v$, where $p_v$ counts positive-charge pairs and $q_v$
counts negative-charge individuals.  The five low charges and their
translated charges are $+,+,+,+,-$.  In the real character field,
orthogonality and the fermion--distinguished STS equations give
\[
 \sum_Z q_Zd_Z\chi_E(FZ)=0,
 \qquad
 \sum_Z\theta_Zq_Zd_Z\chi_E(FZ)=-D\chi_E(g).
\]
Fixed contributions are evaluated from the actual translated restrictions;
for the two-dimensional character one must retain the factor
$\chi_E(g)/\dim(E)=-1$.
The complete 41\,664-assignment integer charge box, independently reduced
to 446 free-variable trials, has 93 solutions on 68 Gram patterns.
The height-six correlated cones keep all 38 individual trace types and
all 28 opposite-pair types.  They also retain the original rotation and
distinguished STS equations.  An actual $Z(\mathrm{Ising})$ control tests
both negative-charge fixed objects and positive-charge pairs, including
all fermion-row normalizations.
The complete correlated campaign excludes 77 allocations by exact dual
certificates and retains 16 allocations, on 16 distinct Gram patterns,
with exact nonnegative rational witnesses.  Thus this strengthened finite
test does not exclude $\sha{R8_3_941}$.
The subsequent supported modular tests at primes $2,3,5,7,11,13$ find no
obstruction.  Bounded integer searches provide twelve nonnegative integer
witnesses, each checked against every original equation; four integer
feasibility questions remain unresolved.  Numerical negative search
outcomes are not used as exclusion certificates.
The additional permutation and mixed-fusion-trace gate leaves fourteen
correlated moment witnesses, eleven with exact integer counts and three
unresolved integer-feasibility cases; its complete finite domain is
described in the accompanying refinement.

The fermionic quotient has four m-type simple orbits.  Its coefficient
tensor has only the trivial and full fusion subrings, so it cannot be a
direct tensor product of two rank-two superfusion categories.  This is a
scoped construction obstruction; no general classification of rank-four
superfusion categories is invoked.  Surviving exact integral or moment
data, wherever recorded in the accompanying receipts, do not supply
super associators or a categorical realization.

\paragraph{A permutation-trace refinement.}
The forced central fermion $G$ in the pseudo-unitary Galois conjugate of
$\sha{R8_3_941}$ supplies another necessary test.  Since $N_G$ is an
involutive permutation matrix and has monodromy eigenvalues $\pm1$, its
number of fixed simples equals $P-N$, where $P$ and $N$ count the two
charges.  Only negative-charge simples can be fixed by balancing.
The ten known distinguished objects and translates form five nonfixed
pairs.  If $p_j$ counts the positive-charge pairs and $q_j$ the remaining
negative-charge high simples of restriction type $j$, then
\[
 P=8+2\sum_jp_j,\qquad N=2+\sum_jq_j,
 \qquad
 \sum_j(q_j\bmod2)\le P-N\le\sum_jq_j,
 \quad P-N\equiv\sum_jq_j\pmod2.
\]
The complete $93$-allocation domain retains $71$ allocations.  Intersecting
with the sixteen exact paired-phase survivors removes allocations $4$
and $86$: they respectively require $-2$ fixed objects, or thirteen
fixed objects among only seven eligible ones.

Products of fusion matrices also have nonnegative integral traces.  The
receipt tests all $420$ quantities
$\operatorname{Tr}(N_G^e\prod_{i=1}^4N_{A_i}^{a_i})$, with
$e\in\{0,1\}$ and $\sum_i a_i\le6$, on every allocation, using the
full simultaneous modular eigenvalues.  An actual $Z(\mathrm{Ising})$
integer-matrix control checks the normalization.  No further survivor is
removed.  The final intersection leaves fourteen rational paired-cone
witnesses, eleven of them with exact nonnegative integer witnesses, and
three unresolved integer-feasibility questions.  These are finite-system
survivals, not constructions of fusion categories.

\subsection{Forced central involutions and correlated semion traces}
\label{subsec:frontier-quartic-semion}

The complete integer Gram and distinguished Verlinde enumerations for
$\sha{R8_3_928}$ and $\sha{R8_3_941}$ each leave $343$ restriction patterns.
Their previously verified Galois--pseudo-unitary gates permit the following
necessary conditions to be used for general complex categorifiability.

\begin{proposition}
\label{prop:frontier-forced-quartic-involution}
Every categorification of either quartic ring, after this Galois transport,
has a unique central invertible $G$ whose restriction is the nontrivial
central source invertible $g$. It satisfies $G^2=\mathbf1$. For
$\sha{R8_3_928}$ one has $\theta_G^2=-1$, whereas for $\sha{R8_3_941}$ one has
$\theta_G^2=1$.
\end{proposition}
\begin{proof}
The complete distinguished restriction list contains only the tensor unit
with restriction supported on $\mathbf1$. Adjunction therefore makes the
forgetful kernel $\mathrm{Vec}$. If a central simple restricts to $m g$,
its product with its dual restricts to $m^2\mathbf1$; the kernel statement
and rigidity force $m=1$. Two such lifts differ by a kernel invertible,
so there is at most one. All $343$ complete integer patterns force exactly
one such lift. The same kernel argument gives $G^2=\mathbf1$.

The complete earlier rotation certificates force
$\nu_2(g)=-1$ and $+1$, respectively. The forgetful functor preserves the
canonical positive pivotal structures. On the pointed order-two
subcategory, the indicator is the normalized cocycle value
$\omega(g,g,g)\in\{1,-1\}$; this also follows from
\cite[Definition~2.2 and Proposition~4.1]{SchauenburgPointedFrontier}
with trivial subgroup. The order-two hexagon gives
$c(G,G)^2=\omega(g,g,g)$, and balancing gives
$1=\theta_G^2c(G,G)^2$. Hence $\theta_G^2=\nu_2(g)$.
\end{proof}

For the semionic case, the monodromy charge of any central simple $Z$
is $\pm1$, since it is an eigenvalue of the involution fusion matrix
$N_G$. Balancing gives
\[
 \theta_{GZ}=\theta_G M_{G,Z}\theta_Z.
\]
Thus $G$ has no fixed simple. The $40$ high restriction vectors fixed by
multiplication by $g$ must have even multiplicity; exactly $15$ of the
$343$ integer patterns remain. The five distinguished low objects
translate to five distinct high objects of multiplicity one, with twists
$\pm i$. Their normalized relative moments at powers $1,\ldots,6$ are
$(0,-1,0,1,0,-1)$, and the second and sixth total moments of every fixed
restriction type vanish.

Imposing these conditions together with the full height-six tensor
rotation equations, all distinguished $STS$ equations, and the fixed
integer Gram pattern gives exact Farkas contradictions for five patterns
and exact rational feasible vectors for ten. This calculation uses the
matrix-character normalization $D/(f_E\dim E)$ throughout.

The remaining correlation can be retained directly. Write
$K=\mathbb Q(\sqrt2,\sqrt3)$ and let $\operatorname{tr}_K$ denote normalized
relative trace. Each unordered $G$-orbit contributes a joint tuple
\[
 \bigl(\operatorname{tr}_K(\theta^k+(i\theta)^k)\bigr)_{k=1}^{6}.
\]
The opposite sign is covered by interchanging the two roots.
Proposition~\ref{prop:m3-biquadratic-traces} gives the finite order box
$288$, together with the all-zero tail. Multiplication by $i$ cannot
remove a prime-power valuation exceeding this bound. Direct exact
averaging of all powers of $\zeta_{288}$ over the $24$ automorphisms fixing
$K$ produces exactly $30$ joint pair types. A pair variable has
multiplicity mass two. Replacing each fixed restriction block by these
pair columns preserves every Gram, rotation, moment, and $STS$ row;
the other high columns and auxiliary rotation variables are retained once.
All ten patterns still have exact rational feasible vectors in the
resulting $873$-equation, $2227$-variable systems.

The simultaneous eigenvalue tuple
\[
 \left(\frac{D\,\chi_E(FZ)}{f_E\dim(E)\dim(Z)}\right)_E
\]
has Galois-invariant multiplicity: the distinguished center fusion
matrices are commuting integer matrices, and Galois transport preserves
the dimensions of their joint eigenspaces. Testing all four automorphisms
of $K$ leaves $255$ of the original $343$ patterns for each quartic ring;
all $15$ semion-parity patterns pass this additional condition.
The same full-tuple test preserves all four complete integer patterns for
$\sha{R8_3_721}$ and all $1238$ cubic/spectral integer patterns for
$\sha{R8_3_1154}$, closing this finite test across the four remaining
noncommutative cases. An actual
$\mathcal Z(\mathrm{Fib})$ control checks the joint-spectrum normalization.

The pair variables themselves count actual orbits and hence are integers;
normalized relative trace changes their values, not their multiplicities.
The auxiliary rotation variables similarly count individual eigenvalues.
A bounded integer search finds exact nonnegative integer witnesses for
four of the ten remaining patterns, verified against all $873$ equations.
The other six integer-feasibility questions remain uncertified: numerical
infeasibility is not used as an obstruction. Exact modular lattice tests
at $2,3,5,7,11,13$, including removal of every variable forced zero by
nonnegative rows, do not resolve them.

These vectors are feasible necessary data, not constructions. In
particular, the forced central semion does not imply a tensor
factorization of the original category. The source-bound read-only
entry point for both finite campaigns is
\begin{quote}\small
\path{r928_semion_rotation/verify_all.py}.
\end{quote}
The independent orbit and correlated-trace checkers are registered
separately. Actual double-semion center data control the $STS$ and twist
normalizations.

\begin{example}[A scope control for fermionic dimension arguments]
\label{ex:frontier-fermion-no-negative-character}
A free central fermion, even in a unitary category and with an all-$M$-type
superquotient, need not admit a scalar ring character taking that fermion
to $-1$.  This supplies a control against a proposed obstruction for
$\sha{R8_3_941}$, rather than a model for that ring.

Let $V=C_2^2$ and
$b(g,h)=g_1h_1+g_2h_2+g_2h_1$ modulo two.  The extension with product
$(g,a)(h,c)=(g+h,a+c+b(g,h))$ is $Q_8$, with central involution
$f=(0,1)$.  In lexicographic order on $\{1,2,3\}^3$, using the encoding
$g=g_1+2g_2$, take the normalized function $A:V^3\to\mathbb Z/4$
whose nonunit table is
\[
 (1,3,0,0,3,1,3,0,1,\;0,3,1,2,1,1,0,0,0,\;3,2,3,0,0,0,3,0,1).
\]
It satisfies
\[
 A(g,h,k)+A(g,h+k,l)+A(h,k,l)
 -A(g+h,k,l)-A(g,h,k+l)=2b(g,h)b(k,l)\pmod4.
\]
Consequently
$\omega((g,a),(h,c),(k,d))=i^{A(g,h,k)+2d b(g,h)}$
is a normalized unitary three-cocycle.  This is the pointed supercocycle
construction of \cite[Theorem~5.4 and Corollary~5.8]{UsherFermionic}.
The receipt also checks all $8^4$ ordinary pentagons directly.

The half-braiding $c_{f,(g,a)}=(-1)^a$ satisfies all central hexagons and
$c_{f,f}=-1$.  Every associator involving two copies of $f$ is one,
so the tensor-square half-braiding is $c_{f,x}^2=1$ and $f\otimes f$
is the unit in the center.  The restriction of $\omega$ to
$\langle f\rangle$ is trivial.  In the positive spherical structure the
twist is $\theta_f=c_{f,f}=-1$; compare
\cite[Section~2.11.5, equation~(26)]{DGNO}.
The underlying $f$-action is free, hence the fermionic quotient has four
$M$-type simple orbits.  Nevertheless $f$ is a commutator in $Q_8$, so
every scalar character of $\mathbb C[Q_8]$ sends $f$ to $+1$.
Thus the absence of a character with value $-1$ is not by itself a
categorification obstruction.  No novelty claim for this pointed
construction is intended.
\end{example}

\begin{proposition}[Complete low-frame convex hull]
Let $R$ be a commutative fusion ring and fix its positive dimension
character. Let $B$ be the induction Gram matrix. For each distinguished
nonunit constituent of $I(\mathbf1)$, let its finite set of possible
restriction columns be fixed by its formal codegree. Enumerate all joint
low frames $L$ with the prescribed unit row, retaining only those for
which $B-LL^{\mathsf T}$ is entrywise nonnegative and positive
semidefinite. Columns with identical allowed sets may be ordered to
remove permutation duplicates. In any finite central-induction linear
relaxation, replace the independent low-column choices by the convex
hull of these complete frames. The resulting system remains necessary
for pseudo-unitary categorification.
\end{proposition}
\begin{proof}
Every actual low restriction matrix has one representative in the
enumeration: the remaining center simples express $B-LL^{\mathsf T}$
as a sum of Gram matrices of nonnegative integral columns. Let $a_c$
denote the full constraint column of a low constituent, including its
Gram and twist-moment contributions. Assign a variable $\lambda_L$ to
each complete frame and use the column $a_L=\sum_{c\in L}a_c$, omitting
the separately fixed tensor unit. The original low-choice equations
become $\sum_L\lambda_L=1$. An actual category gives a one-hot choice
of $L$, so it satisfies the new system. All other columns and the
rotation equations are unchanged.
\end{proof}

The implementation stores these columns implicitly. A row functional
is evaluated on every complete frame by summing its constituent scores;
high columns use an exact factorization into Gram and relative-trace
terms. Column generation proposes row functionals on restricted cones,
but an exclusion is accepted only after checking every column of the
complete finite domain over the integers. Conversely, a recorded
nonnegative rational primal solution is checked against every equation.
Such a solution establishes survival of this necessary-condition
system, not existence of a category.

Applied to $\sha{R8_3_925}$, the complete-frame system at rotation height
$4$ has an exact separating functional for each of the two admissible
second-indicator assignments. The full oriented signed lifting
certificate closes every remaining sphericalization branch, giving an
exclusion over $\mathbb C$ without a pivotality assumption. The
source-bound certificate is
\path{global_induction/frame_results/R8_3_925_h4.json}; its
canonical exact-only replay is
\path{global_induction/verify_accepted_frames.py}.


The same complete-frame criterion at height $6$ excludes
$\sha{R8_3_1739}$. Its domain has $6891$ low frames, $2863$ high
restriction columns, $38$ joint relative-trace types, $323$ equations,
and $115762$ variables. Both admissible second-indicator assignments
have exact full-domain separators. The complete oriented signed lift
also verifies, so this is again an exclusion over $\mathbb C$.
The source-bound certificate is
\path{global_induction/frame_results/R8_3_1739_h6.json},
replayed by the same canonical entrypoint.

\begin{proposition}[Partial modular spectrum and complete character labeling]
Let $R$ be the commutative Grothendieck ring of a pseudo-unitary fusion
category, with global dimension $D$, irreducible characters $\chi_E$,
and formal codegrees $f_E$. If $v$ is the restriction of a simple object
of its center, then
\[
 \frac{D\chi_E(v)}{f_E\dim(v)}
\]
is an algebraic integer for every $E$. Moreover, the distinguished low
restriction columns admit a bijective labeling $v_E$ satisfying
$\dim(v_E)=D/f_E$ and
\[
 (D/f_E)\chi_E(v_F)=(D/f_F)\chi_F(v_E).
\]
\end{proposition}
\begin{proof}
In the unnormalized modular matrix of the center, Deshpande's
induction--restriction character formula, equation~(39) of
\emph{On centers of bimodule categories and induction-restriction
functors} (arXiv:1611.04486), gives
$S_{A_E,Z}=(D/f_E)\chi_E(v_Z)$. Symmetry and the modular character
formula imply that $S_{A_E,Z}/\dim(Z)$ is a fusion-ring character
value, hence integral. Symmetry restricted to the distinguished
objects gives the second assertion.
\end{proof}

The implementation reconstructs and verifies the complete character
table in $\mathbb Q(\zeta_{24})$, whose integer ring is
$\mathbb Z[\zeta_{24}]$, and exhausts all dimension-compatible
bijective low-column labelings. Thus neither filter loses a possible
categorification. For $\sha{R8_3_926}$, the first filter retains $2723$ of
$2863$ high restriction candidates, and symmetric labeling retains
$11$ of $6891$ complete low frames. Exact complete-domain separation
of both indicator cases, followed by all $16$ oriented signed
branches, excludes this ring over $\mathbb C$. The source-bound
certificate is \path{global_induction/spectral_results/R8_3_926_h4.json}.

\begin{proposition}[Distinguished modular symmetry]
Let $R$ be a commutative fusion ring, with positive dimension $D$, ring
characters $\chi_E$, and formal codegrees $f_E$. Enumerate every possible
complete restriction frame of the distinguished central simples in
$I(\mathbf1)$, including all dimension-compatible character labels.
A pseudo-unitary categorification requires at least one such labeling
with
\[
  \frac{D}{f_E}\chi_E(v_F)=\frac{D}{f_F}\chi_F(v_E)
  \qquad\text{for all }E,F.
\]
\end{proposition}
\begin{proof}
The distinguished central simple $A_E$ has dimension $D/f_E$ and
$S(A_E,Z)=(D/f_E)\chi_E(FZ)$\cite[Lemma~4.3 and Corollary~4.5]{DeshpandeCenter}. The assertion is therefore the symmetry
of the unnormalized modular matrix. Every actual frame has a
representative in the complete low-column enumeration, since its
remaining restriction columns give an entrywise nonnegative positive
semidefinite residual Gram matrix. The dimension-compatible bijection
search retains every possible character labeling.
\end{proof}

For $\sha{R8_3_1463}$, the seven raw candidate sets have sizes
$20,9,4,4,12,12,6$. All $2024$ complete frames fail this symmetry test.
The exact character table splits in $\mathbb Q(\sqrt3,\sqrt5)$.
The ordinary character alternatives and all four oriented signed
branches also close, giving an exclusion over $\mathbb C$ without a
pivotality assumption. The canonical exact replay is
\path{global_induction/verify_accepted_generic.py}; no high-column,
phase, STS, or numerical solver computation is required for this proof.

\subsection{A complete octic low-$S$ obstruction for $\sha{R8_3_901}$}
\label{sec:rank8-octic-low-S}
All eight characters of $\sha{R8_3_901}$ lie in
$L=\mathbb Q(\sqrt3,\sqrt5,i)$. The primitive element
$u=\sqrt3+\sqrt5+i$ satisfies
\[
 u^8-28u^6+238u^4+372u^2+441=0.
\]
The exact character adapter checks all multiplication, duality and
orthogonality identities in this field. For the distinguished central
simples, impose
$\beta_E\chi_E(v_F)=\beta_F\chi_F(v_E)$ in all eight rational coordinates,
with no conjugation in this symmetric identity.

The seven raw restriction-column domains have sizes
$6,9,4,4,12,12,20$. The complete induction Gram enumeration leaves
$2024$ low frames. Each has exactly four dimension-compatible character
bijections, and none of the $8096$ bijections satisfies the symmetry
identities. This uses neither high restriction columns nor a
partial-$S$ integrality filter nor linear programming. A separate
checker enumerates full bijections before testing symmetry, independently
of the primary recursive pruning. Its known
$\mathrm{Fib}\boxtimes\mathrm{Vec}_{C_4}$ low-frame control has $576$
permitted bijections, all of which satisfy the identities.

The complete general lift has two ordinary quartic character factors:
the Frobenius--Perron Galois orbit and an orbit with unequal dimensions
on a dual pair. Of the $32$ oriented signed branches, $16$ are gauge
equivalent to the original ring. Each remaining branch has one character
orbit violating duality and one excluded by the norm condition at the
prime $3$, with global-dimension polynomial
\[
 t^4-240t^3+4560t^2-28800t+57600.
\]
Consequently $\sha{R8_3_901}$ admits no complex fusion-category
categorification, without a pivotality or unitarity assumption.
The source-bound replays are
\path{quadratic/octic_low_s_general.py} and
\path{quadratic/review_octic_low_s.py}.

\begin{proposition}[Character-labeled central induction]
Let $R$ be a commutative fusion ring, $K$ its positive dimension field,
and $L$ a verified splitting field of its ring characters. Put
$\beta_E=D/f_E$. In a pseudo-unitary categorification, every central
restriction column $v$ satisfies
\[
 \frac{\beta_E\chi_E(v)}{\dim(v)}\in\mathcal O_L
 \qquad\text{for every }E.
\]
Label each low constituent by its character and impose
$\beta_E\chi_E(v_F)=\beta_F\chi_F(v_E)$ in all rational coordinates
of $L$. These constraints may be imposed either on complete labeled
frames or on independent low-choice simplices, with the latter giving
a weaker necessary relaxation.
\end{proposition}
\begin{proof}
The character formula for central induction gives
$S(A_E,Z)=\beta_E\chi_E(FZ)$
\cite[Lemma~4.3 and Corollary~4.5]{DeshpandeCenter}.
Symmetry of $S$ gives the low equalities, and
$S(A_E,Z)/\dim(Z)$ is a fusion eigenvalue of the center, hence an
algebraic integer. Every actual labeled frame gives one-hot choices
in the independent simplices. Integrality in $L$ is tested by the
integer coefficients of the multiplication characteristic polynomial;
this does not assume that the chosen primitive power basis is an
integral basis.
\end{proof}

If the required character products lie in the relative-trace field,
one may additionally impose the center relation
\[
 \sum_Z\theta_Z\chi_E(v_Z)\chi_F(v_Z)
   =f_F\chi_E(v_F).
\]
This is the unnormalized $STS$ relation, using the center Gauss sum
$D$ and the twist-one distinguished objects. Complete high duality
pairs are combined before taking a relative trace. Their twists agree
and their character coefficients are conjugates, so the pair
coefficient is real. Membership in the actual trace field is checked
exactly before factoring it out of the trace.


\subsection{A centerless rank-five Morita partner excludes \texorpdfstring{$\sha{R8_3_439}$}{R8,3,439}}
\label{subsec:r439-centerless-morita}

The preceding pointed-associator calculation makes full $S_3$
condensation unconditional for $\sha{R8_3_439}$.  The complete dual
based-ring enumeration and frozen rank-five exclusions leave only
$R_{5,3,40}$.  Its basis $(1,s,V,P,Q)$ is self-dual, its dimensions are
$(1,1,2,3+2\sqrt3,3+2\sqrt3)$, and its dimension is
$D=48+24\sqrt3$.  The Galois reduction permits us to take the original
category pseudo-unitary; its Morita dual is then pseudo-unitary by
Morita invariance of categorical and Frobenius--Perron dimensions
\cite[Theorem~2.15 and Corollary~8.14]{ENO};
its canonical spherical structure is provided by
\cite[Proposition~8.23]{ENO}.
Its center has no nontrivial invertibles.  The integral-center
criterion therefore forbids every nonunit integral-dimensional
central simple.

The formal codegrees of this dual ring are $D$, $D'=48-24\sqrt3$,
$3$, $4$, $4$.  The complete induced-unit search gives $24$ low
frames, of which five pass the ordinary moment conditions and just
one passes full distinguished $S$ symmetry:
\[
 (1,0,0,0,0),\quad(1,0,0,1,1),\quad(1,1,1,2,2),\quad
 (1,0,1,1,2),\quad(1,0,1,2,1).
\]
The two assignments of the equal-codegree characters both fix the
simple $A$ corresponding to $D'$.  The remaining Gram matrix is
\[
 Q=\begin{pmatrix}4&2&4&4\\2&6&7&7\\4&7&34&31\\4&7&31&34\end{pmatrix}.
\]
The exact high domain uses positivity of $Q-vv^{\mathsf T}$, the
central simple dimension $d$-number condition
\cite[Theorem~1.8]{OstrikCode}, and algebraic integrality of
$D/d_Z$ \cite[Proposition~3.3]{ENO}.
There are $38$ possible high restriction vectors.  Four have integral
dimension and one violates algebraic integrality of a distinguished
$S$-ratio.  Two independent complete integer enumeration algorithms
find exactly $58$ multiplicity patterns on the remaining vectors.

Let $\chi=(1,1,2,3-2\sqrt3,3-2\sqrt3)$ be the conjugate-FP character.
The distinguished simple $A$ is self-dual: it is the unique simple
of dimension $D/D'=7+4\sqrt3$ among the constituents of the
self-dual object $I(1)$ \cite[Theorem~2.13]{OstrikPiv3}.
The distinguished-row formula \cite[equation~(39)]{DeshpandeCenter}
and the Verlinde formula give
\[
 N_{A,A}^{A}=\frac{D}{(D')^3}
       \sum_{Z\in\operatorname{Irr} Z(\mathcal C)}
       \frac{\chi(FZ)^3}{\operatorname{FPdim}(Z)}.
\]
Counting the five low simples once and each high restriction with its
integer multiplicity gives the following complete list:
\[
\begin{array}{c|ccccc}
N_{A,A}^{A}&(7-\sqrt3)/6&1/3&(19-5\sqrt3)/12&(6-\sqrt3)/3&(1+2\sqrt3)/3\\
\text{number of patterns}&16&16&16&6&4
\end{array}
\]
No entry is an integer.  Thus no possible associator on the forced
rank-five dual can have the required center.  Consequently
$\sha{R8_3_439}$ admits no complex fusion-category realization.

The exact replay reconstructs the complete low-frame domain, both
integer enumerations, and every cubic value; a Fibonacci-center
control checks the cubic normalization.  Its canonical entry point is
\begin{quote}\small\path{evidence/rank8_frontier_20261021/r439_centerless_rank5/verlinde_exclusion.py}\end{quote}
The earlier rational height-six rotation and $STS$ relaxation has a
feasible branch; the integer cubic test supplies the additional
obstruction needed here.

\subsection{A direct cubic central obstruction for \texorpdfstring{$\sha{R8_3_567}$}{R8,3,567}}
\label{subsec:r567-centerless-verlinde}

For $\sha{R8_3_567}$, the general Galois-pseudo-unitary gate and the
integral-center criterion permit the same integer central test directly
on the rank-eight ring.  Its global dimension is $D=48+24\sqrt3$;
its formal codegrees with representation dimensions are
$(D,1)$, $(48-24\sqrt3,1)$, $(12,1)$, $(12,1)$, $(3,2)$.
The full character table is computed in $\mathbb Q(\sqrt3,i)$.
Complete distinguished $S$ symmetry leaves one low frame:
\[
\begin{gathered}
(1,0,0,0,0,0,0,0),\quad(1,0,0,1,1,0,0,0),\\
(1,0,0,0,1,0,0,0),\quad(1,0,0,1,0,0,0,0),\quad
(2,1,1,2,2,0,0,0).
\end{gathered}
\]
The last simple occurs twice in $I(1)$, but once in sums over central
simples. The distinguished row has the normalization
\[
 S_{A_E,Z}=\frac{D}{f_E\dim E}\,\chi_E(FZ)
\]
from \cite[equation~(39)]{DeshpandeCenter} and character orthogonality.

The exact Gram bound, central dimension $d$-number condition
\cite[Theorem~1.8]{OstrikCode}, and algebraic integrality of $D/d_Z$
\cite[Proposition~3.3]{ENO} give $45$ high restriction vectors.
Two have integral dimension and are forbidden.  Nine further vectors
violate algebraic integrality of $S_{A_E,Z}/d_Z$; the complex character
rows are retained in this test.  The remaining $34$ vectors have the
form $(u,u,p,q,w,w,w)$, and their complete Gram equation reduces to
\[
\sum_Z m_Z
\begin{pmatrix}u_Z\\w_Z\\p_Z\\q_Z\end{pmatrix}
\begin{pmatrix}u_Z&w_Z&p_Z&q_Z\end{pmatrix}
=
\begin{pmatrix}
4&0&4&4\\0&4&6&6\\4&6&42&31\\4&6&31&42
\end{pmatrix}.
\]
Two independent complete integer searches agree on $1{,}071$ solutions;
duality of restrictions leaves exactly $155$.

For the self-dual distinguished simple $A$ corresponding to
$D'=48-24\sqrt3$, the cubic coefficient is again
\[
 N_{A,A}^{A}=\frac{D}{(D')^3}\sum_Z\frac{\chi(FZ)^3}{d_Z}.
\]
Its complete list consists of $84$ values $1/3$ and $71$ values $4/3$.
Every integer Gram completion therefore contradicts the Verlinde
formula.  Thus $\sha{R8_3_567}$ has no complex fusion-category realization.
The read-only exact replay, including both enumerations and a known
Fibonacci-center normalization control, is
\begin{quote}\footnotesize\path{evidence/rank8_frontier_20261021/r567_centerless_integer/verlinde_exclusion.py}\end{quote}

\subsection{Integer central data and correlated rotations exclude \texorpdfstring{$\sha{R8_3_455}$}{R8,3,455}}
\label{subsec:r455-integer-rotation}

The general Galois-pseudo-unitary gate also applies to $\sha{R8_3_455}$.
Its dimension field is $\mathbb Q(\sqrt3,\sqrt{13})$, with
$D=39+9\sqrt{13}$.  The complete low/$S$ search has one frame and
$13$ high restriction vectors.  Two independent complete integer Gram
searches give $812$ patterns.  Twenty low cubic Verlinde coefficients
leave $24$; all pass the required field-Galois eigenvalue multiplicity
balance.  When two scalar characters share the FP codegree, the FP
character is explicitly identified as the tensor-unit row.

Two distinguished central simples forget to $\mathbf1$.  Besides the
unit, the other is an invertible $G$ with $G^2=\mathbf1$ and
$\theta_G=1$.  If $S_{G,Z}/d_Z=-1$, balancing gives
$\theta_{GZ}=-\theta_Z$, so $GZ\not\cong Z$.  Since $F(GZ)=F(Z)$,
each such restriction column has even multiplicity.  This leaves
$12$ patterns.  All $120$ low-pair products admit nonnegative integer
aggregate forgetful decompositions; the obstruction requires further
coupling to the twists.

The quadratic conductors are $12,13,156$, and the Cauchy theorem
\cite{BNRWCauchy} supplies twist-order prime support $\{2,3,13\}$.
Proposition~\ref{prop:m3-biquadratic-traces} at height six gives
\[
 B_6=2^4\cdot3^2\cdot13=1872,
\]
and exactly $55$ joint relative-trace types in basis
$(1,\sqrt3,\sqrt{13},\sqrt{39})$.
This bounds the retained trace domain, not the categorical exponent.
The four sign choices in the quadratic Gauss projections are linked
across all retained powers.

For every remaining integer pattern, all four second-indicator choices
are tested against the full divisor-power tensor rotations through
height six, the full distinguished $STS$ equations, and the fixed
integer high multiplicities.  The two equal-codegree assignments give
identical equations because their restriction columns coincide.
Each of the $48$ systems has $381$ equations and $792$ nonnegative
variables.  One common exact integer Farkas vector verifies
$A^{\mathsf T}y\geq0$ and $b^{\mathsf T}y<0$ in every case.
Thus $\sha{R8_3_455}$ has no complex fusion-category realization.

The field adapter is checked against the complete earlier
$\mathbb Q(\sqrt2,\sqrt3)$ trace domain, independent cyclotomic relative
traces, and the actual center $Z(\operatorname{Vec}_{C_2})$ through
height six.  The canonical read-only replay is
\begin{quote}\small\path{evidence/rank8_frontier_20261021/r455_center_integer/verify_all.py}\end{quote}

The same complete labeled-frame $STS$ criterion excludes $\sha{R8_3_903}$.
Its complete low-column enumeration and exact character-label tests leave
four labeled choices. Algebraic integrality retains $2919$ high restriction
vectors, forming $1532$ duality orbits. At height four the necessary system
has $341$ equations and $52127$ columns. Two exact integer functionals
separate all eight permitted indicator assignments. The spherical
alternatives and all $64$ oriented signed branches also close, so the
obstruction holds for general complex fusion categories.
The complete sparse phase-I algorithm is used only for discovery;
\path{global_induction/verify_accepted_complete.py}
reconstructs and verifies the exact whole-domain certificates without
numerical optimization.

The complete labeled-frame $STS$ criterion also excludes $\sha{R8_3_1615}$.
Four labeled low choices survive the complete restriction and
character-label tests. The exact integrality filter retains $3003$
high restriction vectors, each fixed by duality. The height-four
system has $341$ equations and $102141$ columns.
Five exact integer functionals separate all eight permitted indicator
assignments. The spherical alternatives and all eight oriented signed
branches close, giving an obstruction for general complex fusion
categories. The exact-only canonical replay is
\path{global_induction/verify_accepted_complete.py}.

For $\sha{R8_3_1617}$, the complete labeled-frame $STS$ test leaves six
labeled low choices and $3835$ high restriction vectors, each fixed by
duality. The height-four system has $341$ equations and $130431$ columns.
Two exact integer functionals separate the two permitted indicator
assignments. The spherical alternatives and both oriented signed
branches close, proving an exclusion for general complex fusion
categories. The exact-only canonical replay is
\path{global_induction/verify_accepted_complete.py}.

For $\sha{R8_3_1660}$, four labeled low choices and $2937$ high restriction
vectors survive the complete restriction and character tests. All high
vectors are fixed by duality. The height-four $STS$ system has $341$
equations and $99897$ columns. Three exact integer functionals separate
all eight permitted indicator assignments. The spherical alternatives
and all eight oriented signed branches close, proving an exclusion for
general complex fusion categories. The exact-only canonical replay is
\path{global_induction/verify_accepted_complete.py}.

For $\sha{R8_3_1661}$, the complete labeled-frame $STS$ test leaves four
labeled low choices and $3051$ high restriction vectors, each fixed by
duality. The height-four system has $341$ equations and $103773$ columns.
Two exact integer functionals separate the two permitted indicator
assignments. The spherical alternatives and all six oriented signed
branches close, proving an exclusion for general complex fusion
categories. The exact-only canonical replay is
\path{global_induction/verify_accepted_complete.py}.

This complete labeled-frame system at height $4$ excludes
$\sha{R8_3_985}$. The modular filters leave four labeled low frames from
$118696$ complete low frames and $3045$ high restriction vectors.
Both indicator assignments have exact full-domain separating
functionals; all $24$ oriented signed branches also close.
The canonical replay is
\path{global_induction/verify_accepted_sts.py}.

The full-field atomic symmetry system excludes $\sha{R8_3_3413}$ at
height $4$. Its characters split in the real degree-six field
$\mathbb Q(\sqrt5,\zeta_9+\zeta_9^{-1})$. The cone has $253$ equations,
$4485$ variables, $367$ high restriction vectors and ten joint trace
types. Its sole indicator assignment has an exact separating
functional, and its complete signed lift verifies. This proof uses
full-field symmetry and integrality, without factoring character
products outside the phase-trace field. The canonical replay is
\path{global_induction/verify_accepted_atomic.py}.

The full-field atomic spectral criterion also excludes $\sha{R8_3_1055}$.
Its characters split in the field
\[
  \mathbb Q(\sqrt{37},\sqrt{-3}).
\]
Exact algebraic integrality leaves
$566$ high restriction vectors in $352$ duality orbits. The
height-four cone has $211$ equations and $16695$ columns; its full
complex $S$-symmetry equations do not require a phase-trace extension.
Both indicator assignments have exact separating functionals.
The spherical alternatives and all four oriented signed branches
also close, giving a general complex exclusion. Independent normal
and optimized replays check every column for both rays.
The canonical exact-only replay is
\path{global_induction/verify_accepted_atomic.py}.

The character-labeled atomic $STS$ relaxation also excludes
$\sha{R8_3_3258}$. All eight characters split in its real dimension field
$\mathbb Q(\sqrt5)$. The exact spectral filter leaves $1471$ high
restriction vectors, and the height-four cone has $241$ equations,
$17619$ variables, and ten joint phase-trace types. Its sole indicator
assignment has an exact full-domain separating functional. The
spherical alternatives and the complete oriented signed lift also
verify, so this is an exclusion over $\mathbb C$.
The canonical exact-only replay is
\path{global_induction/verify_accepted_atomic_sts.py}.
